\moisgauche{Décembre}{2026}{\devmois{Mois de l’Immaculée Conception}{364}{1}}{Cieux, répandez d'en-haut votre rosée, et que les nuées fassent pleuvoir le Juste ; que la terre s'ouvre, et qu'elle fasse germer le Sauveur.}{Isa 45:8}{4\textsuperscript{e} dimanche de l’Avent}
\moisdroite{Décembre 2026}{\rubmois{Fêtes principales}\lignemois{3}{S. François-Xavier, Confesseur}\lignemois{8}{Immaculée Conception de la Bienheureuse Vierge Marie}\lignemois{15}{Octave de l’Immaculée Conception}\lignemois{21}{S. Thomas, Apôtre}\lignemois{25}{Nativité de Notre-Seigneur Jésus-Christ}\lignemois{26}{S. Étienne, Protomartyr}\lignemois{27}{S. Jean, Apôtre et Évangéliste}\lignemois{28}{Les Saints Innocents}\rubmois{Dévotions et neuvaines qui commencent}\lignemois{4}{Premier vendredi du mois\pg{69}}\lignemois{5}{Premier samedi du mois\pg{73}}\lignemois{16}{Neuvaine de Noël}\lignemois{20}{Six dimanches de S. Thomas d’Aquin\pg{64}}\rubmois{Pénitence}\lignemoislong{Jeûne et abstinence}{16, 18, 19, 24}\lignemoislong{Abstinence}{4, 11}}
\semainegauche{du 29 novembre au 5 décembre 2026}{1\textsuperscript{er} dimanche de l’Avent}
\jourcal{Dimanche}{29}{1}{violet}{\offgrand{1\textsuperscript{er} dimanche de l’Avent}}{dim. 1\textsuperscript{re} cl.}{\comm{S. Saturnin, Martyr}\liturgie{Intr. \textit{Ad te levávi ánimam meam} \enspace·\enspace \textit{Statio ad S. Mariam majorem}}\nov{Neuvaine à l’Immaculée Conception — 1\textsuperscript{er} jour}{}}
\jourfete{Lundi}{30}{1}{rouge}{\offgrand{S. André, Apôtre}}{2\textsuperscript{e} cl.}{\comm{Férie de l’Avent}\liturgie{Intr. \textit{Mihi autem nimis honoráti sunt…}}\noticesaint{À Patras, en Achaïe, l’anniversaire de l’Apôtre saint André, qui prêcha le saint Évangile du Christ dans la Thrace et dans la Scythie. […]}\nov{Neuvaine à l’Immaculée Conception — 2\textsuperscript{e} jour}{}}
\jourcal{Mardi}{1}{0}{violet}{\offpetit{Férie de l’Avent}}{fér. maj.}{\liturgie{Intr. \textit{Ad te levávi ánimam meam}}\nov{Neuvaine à l’Immaculée Conception — 3\textsuperscript{e} jour}{}}
\jourcal{Mercredi}{2}{0}{rouge}{\offpetit{Ste Bibiane, Vierge et Martyre}}{semi-dbl}{\comm{Férie de l’Avent}\liturgie{Intr. \textit{Me exspectavérunt peccatóres}}\dev{Premier mercredi : S. Joseph}{\pg{68}}\nov{Neuvaine à l’Immaculée Conception — 4\textsuperscript{e} jour}{}}
\semainedroite{du 29 novembre au 5 décembre 2026}
\jourcal{Jeudi}{3}{0}{blanc}{\offgrand{S. François-Xavier, Confesseur}}{dbl maj.}{\comm{Férie de l’Avent}\liturgie{Intr. \textit{Loquébar de testimóniis tuis in…}}\noticesaint{Dans l'île de San-Chan, dans la rivière Canton, en Chine, saint François Xavier de la Compagnie de Jésus, apôtre des Indes, illustre pour la conversion des païens, les dons et les miracles, qui est décédé (en […]}\dev{Heure sainte (veille du premier vendredi)}{\pg{69}}\nov{Neuvaine à l’Immaculée Conception — 5\textsuperscript{e} jour}{}}
\jourcal{Vendredi}{4}{0}{blanc}{\offpetit{S. Pierre Chrysologue, Évêque, Confesseur et Docteur de l’Église}}{dbl}{\comm{Férie de l’Avent ; Ste Barbe, Vierge et Martyre}\liturgie{Intr. \textit{In médio Ecclésiæ apéruit os…}}\jeune{Abstinence}\dev{Premier vendredi du mois — 1/9}{\pg{69}}\nov{Neuvaine à l’Immaculée Conception — 6\textsuperscript{e} jour}{}}
\jourcal{Samedi}{5}{0}{violet}{\offpetit{Férie de l’Avent}}{fér. maj.}{\comm{S. Sabba, Abbé}\liturgie{Intr. \textit{Ad te levávi ánimam meam}}\dev{Premier samedi du mois — 1/5}{\pg{73}}\nov{Neuvaine à l’Immaculée Conception — 7\textsuperscript{e} jour}{}}
\notes
\semainegauche{du 6 au 12 décembre 2026}{2\textsuperscript{e} dimanche de l’Avent}
\jourcal{Dimanche}{6}{1}{violet}{\offgrand{2\textsuperscript{e} dimanche de l’Avent}}{dim. 2\textsuperscript{e} cl.}{\comm{S. Nicolas, Évêque et Confesseur}\liturgie{Intr. \textit{Pópulus Sion} \enspace·\enspace \textit{Statio ad S. Crucem in Jerusalem}}\nov{Neuvaine à l’Immaculée Conception — 8\textsuperscript{e} jour}{}}
\jourcal{Lundi}{7}{0}{blanc}{\offpetit{S. Ambroise, Évêque, Confesseur et Docteur de l’Église}}{dbl}{\comm{Férie de l’Avent}\liturgie{Intr. \textit{In médio Ecclésiæ apéruit os…}}\nov{Neuvaine à l’Immaculée Conception — 9\textsuperscript{e} jour}{}}
\jourfete{Mardi}{8}{1}{blanc}{\offgrand{Immaculée Conception de la Bienheureuse Vierge Marie}}{1\textsuperscript{re} cl.}{\comm{Férie de l’Avent}\liturgie{Intr. \textit{Gaudens gaudébo in Dómino}}\noticesaint{La Conception Immaculée de la glorieuse et toujours Vierge Marie, Mère de Dieu, qui, par un insigne privilège de Dieu, fut, comme l’a solennellement défini en ce jour le souverain pontife Pie IX, totalement préservée de la tache du péché originel.}}
\jourcal{Mercredi}{9}{0}{blanc}{\offpetit{2\textsuperscript{e} jour dans l’octave de l’Immaculée Conception}}{semi-dbl}{\comm{Férie de l’Avent}\liturgie{Intr. \textit{Gaudens gaudébo in Dómino}}}
\semainedroite{du 6 au 12 décembre 2026}
\jourcal{Jeudi}{10}{0}{blanc}{\offpetit{3\textsuperscript{e} jour dans l’octave de l’Immaculée Conception}}{semi-dbl}{\comm{Férie de l’Avent ; S. Melchiade, Pape et Martyr}\liturgie{Intr. \textit{Gaudens gaudébo in Dómino}}}
\jourcal{Vendredi}{11}{0}{blanc}{\offpetit{S. Damase I\textsuperscript{er}, Pape et Confesseur}}{semi-dbl}{\comm{Férie de l’Avent ; 4\textsuperscript{e} jour dans l’octave de l’Immaculée Conception}\liturgie{Intr. \textit{Státuit ei Dóminus testaméntum pacis}}\jeune{Abstinence}}
\jourcal{Samedi}{12}{0}{blanc}{\offpetit{5\textsuperscript{e} jour dans l’octave de l’Immaculée Conception}}{semi-dbl}{\comm{Férie de l’Avent}\liturgie{Intr. \textit{Gaudens gaudébo in Dómino}}}
\notes
\semainegauche{du 13 au 19 décembre 2026}{3\textsuperscript{e} dimanche de l’Avent (Gaudete)}
\jourcal{Dimanche}{13}{1}{rose}{\offgrand{3\textsuperscript{e} dimanche de l’Avent (Gaudete)}}{dim. 2\textsuperscript{e} cl.}{\comm{Ste Lucie, Vierge et Martyre ; 6\textsuperscript{e} jour dans l’octave de l’Immaculée Conception}\liturgie{Intr. \textit{Gaudéte in Dómino semper} \enspace·\enspace \textit{Statio ad S. Petrum}}}
\jourcal{Lundi}{14}{0}{blanc}{\offpetit{7\textsuperscript{e} jour dans l’octave de l’Immaculée Conception}}{semi-dbl}{\comm{Férie de l’Avent}\liturgie{Intr. \textit{Gaudens gaudébo in Dómino}}}
\jourcal{Mardi}{15}{0}{blanc}{\offgrand{Octave de l’Immaculée Conception}}{dbl maj.}{\comm{Férie de l’Avent}\liturgie{Intr. \textit{Gaudens gaudébo in Dómino}}}
\jourcal{Mercredi}{16}{0}{rouge}{\offpetit{S. Eusèbe, Évêque et Martyr}}{semi-dbl}{\comm{Mercredi des Quatre-Temps de l’Avent}\liturgie{Intr. \textit{Sacerdótes Dei} \enspace·\enspace \textit{Statio ad S. Mariam majorem}}\jeune{Jeûne et abstinence}\nov{Neuvaine de Noël — 1\textsuperscript{er} jour}{}}
\semainedroite{du 13 au 19 décembre 2026}
\jourcal{Jeudi}{17}{0}{violet}{\offpetit{Férie de l’Avent}}{fér. maj.}{\liturgie{Intr. \textit{Gaudéte in Dómino semper}}\nov{Neuvaine de Noël — 2\textsuperscript{e} jour}{}}
\jourcal{Vendredi}{18}{0}{violet}{\offpetit{Vendredi des Quatre-Temps de l’Avent}}{fér. maj.}{\liturgie{Intr. \textit{Prope es tu} \enspace·\enspace \textit{Statio ad Ss. duodecim Apostolos}}\jeune{Jeûne et abstinence}\nov{Neuvaine de Noël — 3\textsuperscript{e} jour}{}}
\jourcal{Samedi}{19}{0}{violet}{\offpetit{Samedi des Quatre-Temps de l’Avent}}{fér. maj.}{\comm{Vigile de S. Thomas, Apôtre}\liturgie{Intr. \textit{Veni} \enspace·\enspace \textit{Statio ad S. Petrum}}\jeune{Jeûne et abstinence}\nov{Neuvaine de Noël — 4\textsuperscript{e} jour}{}}
\notes
\semainegauche{du 20 au 26 décembre 2026}{4\textsuperscript{e} dimanche de l’Avent}
\jourcal{Dimanche}{20}{1}{violet}{\offgrand{4\textsuperscript{e} dimanche de l’Avent}}{dim. 2\textsuperscript{e} cl.}{\liturgie{Intr. \textit{Roráte} \enspace·\enspace \textit{Statio ad Ss. duodecim Apostolos}}\dev{Six dimanches de S. Thomas d’Aquin — 1/6}{\pg{64}}\nov{Neuvaine de Noël — 5\textsuperscript{e} jour}{}}
\jourfete{Lundi}{21}{1}{rouge}{\offgrand{S. Thomas, Apôtre}}{2\textsuperscript{e} cl.}{\comm{Férie de l’Avent}\liturgie{Intr. \textit{Mihi autem nimis honoráti sunt…}}\noticesaint{À Calamine, l’anniversaire du bienheureux Apôtre Thomas, qui prêcha l’Évangile aux Parthes, aux Mèdes, aux Perses et aux Hyrcaniens. […]}\nov{Neuvaine de Noël — 6\textsuperscript{e} jour}{}}
\jourcal{Mardi}{22}{0}{violet}{\offpetit{Férie de l’Avent}}{fér. maj.}{\liturgie{Intr. \textit{Roráte}}\nov{Neuvaine de Noël — 7\textsuperscript{e} jour}{}}
\jourcal{Mercredi}{23}{0}{violet}{\offpetit{Férie de l’Avent}}{fér. maj.}{\liturgie{Intr. \textit{Roráte}}\nov{Neuvaine de Noël — 8\textsuperscript{e} jour}{}}
\semainedroite{du 20 au 26 décembre 2026}
\jourcal{Jeudi}{24}{0}{violet}{\offpetit{Vigile de la Nativité de Notre-Seigneur}}{vigile 1\textsuperscript{re} cl.}{\liturgie{Intr. \textit{Hódie sciétis}}\jeune{Jeûne et abstinence}\nov{Neuvaine de Noël — 9\textsuperscript{e} jour}{}}
\jourfete{Vendredi}{25}{1}{blanc}{\offgrand{Nativité de Notre-Seigneur Jésus-Christ}}{1\textsuperscript{re} cl.}{\liturgie{Intr. \textit{Puer natus est nobis} \enspace·\enspace \textit{Statio ad S. Mariam majorem ad Praesepe ; ad S. Anastasiam ; ad S. Mariam majorem}}\noticesaint{Depuis la création du monde, quand Dieu au commencement tira du néant le ciel et la terre, l’an cinq mille cent quatre-vingt-dix-neuf ; depuis le déluge, l’an deux mille neuf cent cinquante-sept ; […]}}
\jourfete{Samedi}{26}{1}{rouge}{\offgrand{S. Étienne, Protomartyr}}{2\textsuperscript{e} cl.}{\liturgie{Intr. \textit{Sedérunt príncipes} \enspace·\enspace \textit{Statio ad S. Stephanum in Cælio monte}}\noticesaint{À Jérusalem, l’anniversaire de saint Étienne, premier Martyr, qui fut lapidé par les Juifs, peu de temps après l’Ascension du Seigneur.}}
\notes
\moisgauche{Janvier}{2027}{\devmois{Mois du saint Nom de Jésus}{118}{1}}{Voilà que vient le Seigneur Maître ; le pouvoir est dans sa main, la puissance et l’empire.}{Malach 3:1; 1 Par 29:12}{Épiphanie}
\moisdroite{Janvier 2027}{\rubmois{Fêtes principales}\lignemois{1}{Circoncision de Notre-Seigneur et octave de Noël}\lignemois{3}{Très Saint Nom de Jésus}\lignemois{6}{Épiphanie de Notre-Seigneur}\lignemois{13}{Octave de l’Épiphanie (Baptême du Seigneur)}\lignemois{18}{Chaire de S. Pierre à Rome}\lignemois{25}{Conversion de S. Paul, Apôtre}\rubmois{Dévotions et neuvaines qui commencent}\lignemois{8}{Treize vendredis de S. François de Paule\pg{72}}\lignemois{31}{Sept dimanches de S. Joseph\pg{64}}\rubmois{Pénitence}\lignemoislong{Abstinence}{1, 8, 15, 22, 29}}
\semainegauche{du 27 décembre 2026 au 2 janvier 2027}{S. Jean, Apôtre et Évangéliste}
\jourfete{Dimanche}{27}{1}{blanc}{\offgrand{S. Jean, Apôtre et Évangéliste}}{2\textsuperscript{e} cl.}{\liturgie{Intr. \textit{In médio Ecclésiæ apéruit os…} \enspace·\enspace \textit{Statio ad S. Mariam majorem}}\noticesaint{À Éphèse, l’anniversaire de saint Jean, Apôtre et Évangéliste. […]}\dev{Six dimanches de S. Thomas d’Aquin — 2/6}{\pg{64}}}
\jourfete{Lundi}{28}{1}{violet}{\offgrand{Les Saints Innocents}}{2\textsuperscript{e} cl.}{\liturgie{Intr. \textit{Ex ore infántium} \enspace·\enspace \textit{Statio ad S. Paulum}}\noticesaint{À Bethléem de Juda, l’anniversaire des saints Innocents, qui pour le Christ furent massacrés par le roi Hérode.}}
\jourcal{Mardi}{29}{0}{rouge}{\offpetit{S. Thomas, Évêque et Martyr}}{dbl}{\comm{Octave de Noël}\liturgie{Intr. \textit{Gaudeámus omnes in Dómino}}}
\jourcal{Mercredi}{30}{0}{blanc}{\offpetit{Dimanche dans l’octave de Noël}}{dim.}{\liturgie{Intr. \textit{Dum médium siléntium tenérent ómnia}}}
\semainedroite{du 27 décembre 2026 au 2 janvier 2027}
\jourcal{Jeudi}{31}{0}{blanc}{\offpetit{S. Sylvestre, Pape et Confesseur}}{dbl}{\comm{Octave de Noël}\liturgie{Intr. \textit{Státuit ei Dóminus testaméntum pacis}}\dev{Heure sainte (veille du premier vendredi)}{\pg{69}}}
\jourfete{Vendredi}{1}{1}{blanc}{\offgrand{Circoncision de Notre-Seigneur et octave de Noël}}{2\textsuperscript{e} cl.}{\liturgie{Intr. \textit{Puer natus est nobis} \enspace·\enspace \textit{Statio ad S. Mariam trans Tiberim}}\jeune{Abstinence (sauf là où la Circoncision est fête de précepte)}\dev{Premier vendredi du mois — 2/9}{\pg{69}}}
\jourcal{Samedi}{2}{0}{rouge}{\offpetit{Octave de saint Étienne}}{simple}{\liturgie{Intr. \textit{Sedérunt príncipes}}\dev{Premier samedi du mois — 2/5}{\pg{73}}}
\notes
\semainegauche{du 3 au 9 janvier 2027}{Très Saint Nom de Jésus}
\jourcal{Dimanche}{3}{1}{blanc}{\offgrand{Très Saint Nom de Jésus}}{2\textsuperscript{e} cl.}{\liturgie{Intr. \textit{In nómine Jesu omne genu…}}\dev{Six dimanches de S. Thomas d’Aquin — 3/6}{\pg{64}}}
\jourcal{Lundi}{4}{0}{rouge}{\offpetit{Octave des Saints Innocents}}{simple}{\liturgie{Intr. \textit{Ex ore infántium}}}
\jourcal{Mardi}{5}{0}{blanc}{\offpetit{Vigile de l’Épiphanie}}{vigile}{\comm{S. Télesphore, Pape et Martyr}}
\jourfete{Mercredi}{6}{1}{blanc}{\offgrand{Épiphanie de Notre-Seigneur}}{1\textsuperscript{re} cl.}{\liturgie{Intr. \textit{Ecce} \enspace·\enspace \textit{Statio ad S. Petrum}}\noticesaint{L’Épiphanie du Seigneur.}\dev{Premier mercredi : S. Joseph}{\pg{68}}}
\semainedroite{du 3 au 9 janvier 2027}
\jourcal{Jeudi}{7}{0}{blanc}{\offpetit{2\textsuperscript{e} jour dans l’octave de l’Épiphanie}}{semi-dbl}{\liturgie{Intr. \textit{Ecce}}}
\jourcal{Vendredi}{8}{0}{blanc}{\offpetit{3\textsuperscript{e} jour dans l’octave de l’Épiphanie}}{semi-dbl}{\liturgie{Intr. \textit{Ecce}}\jeune{Abstinence}\dev{Treize vendredis de S. François de Paule — 1/13}{\pg{72}}}
\jourcal{Samedi}{9}{0}{blanc}{\offpetit{4\textsuperscript{e} jour dans l’octave de l’Épiphanie}}{semi-dbl}{\liturgie{Intr. \textit{Ecce}}}
\notes
\semainegauche{du 10 au 16 janvier 2027}{Sainte Famille de Jésus, Marie et Joseph}
\jourcal{Dimanche}{10}{1}{blanc}{\offgrand{Sainte Famille de Jésus, Marie et Joseph}}{dbl maj.}{\comm{5\textsuperscript{e} jour dans l’octave de l’Épiphanie}\liturgie{Intr. \textit{Exsúltet gáudio pater Justi}}\dev{Six dimanches de S. Thomas d’Aquin — 4/6}{\pg{64}}}
\jourcal{Lundi}{11}{0}{blanc}{\offpetit{6\textsuperscript{e} jour dans l’octave de l’Épiphanie}}{semi-dbl}{\comm{S. Hygin, Pape et Martyr}\liturgie{Intr. \textit{Státuit ei Dóminus testaméntum pacis}}}
\jourcal{Mardi}{12}{0}{blanc}{\offpetit{7\textsuperscript{e} jour dans l’octave de l’Épiphanie}}{semi-dbl}{\liturgie{Intr. \textit{Ecce}}}
\jourcal{Mercredi}{13}{0}{blanc}{\offgrand{Octave de l’Épiphanie (Baptême du Seigneur)}}{dbl maj.}{\liturgie{Intr. \textit{Ecce}}}
\semainedroite{du 10 au 16 janvier 2027}
\jourcal{Jeudi}{14}{0}{blanc}{\offpetit{S. Hilaire, Évêque, Confesseur et Docteur de l’Église}}{dbl}{\comm{S. Félix, Prêtre et Martyr}\liturgie{Intr. \textit{In médio Ecclésiæ apéruit os…}}}
\jourcal{Vendredi}{15}{0}{blanc}{\offpetit{S. Paul, Premier Ermite, Confesseur}}{dbl}{\comm{S. Maur, Abbé}\liturgie{Intr. \textit{Justus ut palma florébit}}\jeune{Abstinence}\dev{Treize vendredis de S. François de Paule — 2/13}{\pg{72}}}
\jourcal{Samedi}{16}{0}{rouge}{\offpetit{S. Marcel, Pape et Martyr}}{semi-dbl}{\liturgie{Intr. \textit{Státuit ei Dóminus testaméntum pacis}}}
\notes
\semainegauche{du 17 au 23 janvier 2027}{2\textsuperscript{e} dimanche après l’Épiphanie}
\jourcal{Dimanche}{17}{1}{vert}{\offgrand{2\textsuperscript{e} dimanche après l’Épiphanie}}{dim.}{\comm{S. Antoine, Abbé}\liturgie{Intr. \textit{Omnis terra adóret te}}\dev{Six dimanches de S. Thomas d’Aquin — 5/6}{\pg{64}}}
\jourcal{Lundi}{18}{0}{blanc}{\offgrand{Chaire de S. Pierre à Rome}}{dbl maj.}{\comm{Ste Prisque, Vierge et Martyre}\liturgie{Intr. \textit{Státuit ei Dóminus testaméntum pacis}}\noticesaint{La Chaire de saint Pierre, Apôtre, sur laquelle il siégea le premier, à Rome.}}
\jourcal{Mardi}{19}{0}{rouge}{\offpetit{Ss. Marius, Marthe, Audifax et Abacum, Martyrs}}{simple}{\comm{S. Canut, Roi et Martyr}\liturgie{Intr. \textit{Justi epuléntur et exsúltent in…}}}
\jourcal{Mercredi}{20}{0}{rouge}{\offpetit{Ss. Fabien et Sébastien, Martyrs}}{dbl}{\liturgie{Intr. \textit{Intret in conspéctu tuo}}}
\semainedroite{du 17 au 23 janvier 2027}
\jourcal{Jeudi}{21}{0}{rouge}{\offpetit{Ste Agnès, Vierge et Martyre}}{dbl}{\liturgie{Intr. \textit{Me exspectavérunt peccatóres}}}
\jourcal{Vendredi}{22}{0}{rouge}{\offpetit{Ss. Vincent et Anastase, Martyrs}}{semi-dbl}{\liturgie{Intr. \textit{Intret in conspéctu tuo}}\jeune{Abstinence}\dev{Treize vendredis de S. François de Paule — 3/13}{\pg{72}}}
\jourcal{Samedi}{23}{0}{vert}{\offpetit{3\textsuperscript{e} dimanche après l’Épiphanie (anticipé au samedi)}}{dim.}{\liturgie{Intr. \textit{Adoráte Deum}}}
\notes
\semainegauche{du 24 au 30 janvier 2027}{Dimanche de la Septuagésime}
\jourcal{Dimanche}{24}{1}{violet}{\offgrand{Dimanche de la Septuagésime}}{dim. 2\textsuperscript{e} cl.}{\comm{S. Timothée, Évêque et Martyr}\liturgie{Intr. \textit{Circumdedérunt me gémitus mortis} \enspace·\enspace \textit{Statio ad S. Laurentium extra muros}}\dev{Six dimanches de S. Thomas d’Aquin — 6/6}{\pg{64}}}
\jourcal{Lundi}{25}{0}{blanc}{\offgrand{Conversion de S. Paul, Apôtre}}{dbl maj.}{\liturgie{Intr. \textit{Scio}}\noticesaint{La Conversion de l’Apôtre saint Paul, qui eut lieu la seconde année après l’Ascension du Seigneur.}}
\jourcal{Mardi}{26}{0}{rouge}{\offpetit{S. Polycarpe, Évêque et Martyr}}{dbl}{\liturgie{Intr. \textit{Sacerdótes Dei}}}
\jourcal{Mercredi}{27}{0}{blanc}{\offpetit{S. Jean Chrysostome, Évêque, Confesseur et Docteur de l’Église}}{dbl}{\liturgie{Intr. \textit{In médio Ecclésiæ apéruit os…}}}
\semainedroite{du 24 au 30 janvier 2027}
\jourcal{Jeudi}{28}{0}{blanc}{\offpetit{S. Pierre Nolasque, Confesseur}}{dbl}{\comm{Ste Agnès, Vierge et Martyre, pour la seconde fois}\liturgie{Intr. \textit{Justus ut palma florébit}}}
\jourcal{Vendredi}{29}{0}{blanc}{\offpetit{S. François de Sales, Évêque, Confesseur et Docteur de l’Église}}{dbl}{\liturgie{Intr. \textit{In médio Ecclésiæ apéruit os…}}\jeune{Abstinence}\dev{Treize vendredis de S. François de Paule — 4/13}{\pg{72}}}
\jourcal{Samedi}{30}{0}{rouge}{\offpetit{Ste Martine, Vierge et Martyre}}{semi-dbl}{\liturgie{Intr. \textit{Loquébar de testimóniis tuis in…}}}
\notes
\moisgauche{Février}{2027}{}{Nous avons reçu, ô Dieu, votre miséricorde au milieu de votre temple : comme votre nom, ô Dieu, ainsi votre louange s’étend jusqu’aux extrémités de la terre : votre droite est pleine de justice.}{Ps 47:10-11}{Purification}
\moisdroite{Février 2027}{\rubmois{Fêtes principales}\lignemois{2}{Purification de la Bienheureuse Vierge Marie}\lignemois{11}{Apparition de la Bienheureuse Vierge Marie Immaculée (Lourdes)}\lignemois{22}{Chaire de S. Pierre, Apôtre}\lignemois{24}{S. Matthias, Apôtre}\rubmois{Pénitence}\lignemoislong{Jeûne et abstinence}{10, 12, 13, 17, 19, 20, 26, 27}\lignemoislong{Jeûne}{11, 15, 16, 18, 22, 23, 24, 25}\lignemoislong{Abstinence}{5}}
\semainegauche{du 31 janvier au 6 février 2027}{Dimanche de la Sexagésime}
\jourcal{Dimanche}{31}{1}{violet}{\offgrand{Dimanche de la Sexagésime}}{dim. 2\textsuperscript{e} cl.}{\comm{S. Jean Bosco, Confesseur}\liturgie{Intr. \textit{Exsúrge} \enspace·\enspace \textit{Statio ad S. Paulum}}\dev{Sept dimanches de S. Joseph — 1/7}{\pg{64}}}
\jourcal{Lundi}{1}{0}{rouge}{\offpetit{S. Ignace d’Antioche, Évêque et Martyr}}{dbl}{\liturgie{Intr. \textit{Mihi autem absit gloriári}}}
\jourfete{Mardi}{2}{1}{blanc}{\offgrand{Purification de la Bienheureuse Vierge Marie}}{2\textsuperscript{e} cl.}{\liturgie{Intr. \textit{Suscépimus}}\noticesaint{La Purification de la Bienheureuse Vierge Marie, fête appelée par les Grecs Hypapante (ou Rencontre) du Seigneur.}}
\jourcal{Mercredi}{3}{0}{rouge}{\offpetit{S. Blaise, Évêque et Martyr}}{simple}{\liturgie{Intr. \textit{Sacerdótes Dei}}\dev{Premier mercredi : S. Joseph}{\pg{68}}}
\semainedroite{du 31 janvier au 6 février 2027}
\jourcal{Jeudi}{4}{0}{blanc}{\offpetit{S. André Corsini, Évêque et Confesseur}}{dbl}{\liturgie{Intr. \textit{Státuit ei Dóminus testaméntum pacis}}\dev{Heure sainte (veille du premier vendredi)}{\pg{69}}}
\jourcal{Vendredi}{5}{0}{rouge}{\offpetit{Ste Agathe, Vierge et Martyre}}{dbl}{\liturgie{Intr. \textit{Gaudeámus omnes in Dómino}}\jeune{Abstinence}\dev{Treize vendredis de S. François de Paule — 5/13}{\pg{72}}\dev{Premier vendredi du mois — 3/9}{\pg{69}}}
\jourcal{Samedi}{6}{0}{blanc}{\offpetit{S. Tite, Évêque et Confesseur}}{dbl}{\comm{Ste Dorothée, Vierge et Martyre}\liturgie{Intr. \textit{Státuit ei Dóminus testaméntum pacis}}\dev{Premier samedi du mois — 3/5}{\pg{73}}}
\notes
\semainegauche{du 7 au 13 février 2027}{Dimanche de la Quinquagésime}
\jourcal{Dimanche}{7}{1}{violet}{\offgrand{Dimanche de la Quinquagésime}}{dim. 2\textsuperscript{e} cl.}{\comm{S. Romuald, Abbé}\liturgie{Intr. \textit{Esto mihi in Deum protectórem} \enspace·\enspace \textit{Statio ad S. Petrum}}\dev{Sept dimanches de S. Joseph — 2/7}{\pg{64}}}
\jourcal{Lundi}{8}{0}{blanc}{\offpetit{S. Jean de Matha, Confesseur}}{dbl}{\liturgie{Intr. \textit{Os justi meditábitur sapiéntiam}}}
\jourcal{Mardi}{9}{0}{blanc}{\offpetit{S. Cyrille d’Alexandrie, Évêque, Confesseur et Docteur de l’Église}}{dbl}{\comm{Ste Apollonie, Vierge et Martyre}\liturgie{Intr. \textit{In médio Ecclésiæ apéruit os…}}}
\jourcal{Mercredi}{10}{0}{violet}{\offpetit{Mercredi des Cendres}}{fér. priv.}{\comm{Ste Scholastique, Vierge}\liturgie{Intr. \textit{Miseréris ómnium} \enspace·\enspace \textit{Statio ad S. Sabinam}}\jeune{Jeûne et abstinence}}
\semainedroite{du 7 au 13 février 2027}
\jourcal{Jeudi}{11}{0}{blanc}{\offgrand{Apparition de la Bienheureuse Vierge Marie Immaculée (Lourdes)}}{dbl maj.}{\comm{Jeudi après les Cendres}\liturgie{Intr. \textit{Vidi civitátem sanctam} \enspace·\enspace \textit{Statio ad S. Georgium}}\noticesaint{À Lourdes, en France, l’apparition de la Bienheureuse Vierge Marie Immaculée.}\jeune{Jeûne}}
\jourcal{Vendredi}{12}{0}{blanc}{\offpetit{Sept Saints Fondateurs de l’Ordre des Servites de la Bse Vierge Marie}}{dbl}{\comm{Vendredi après les Cendres}\liturgie{Intr. \textit{Justi decantavérunt} \enspace·\enspace \textit{Statio ad Ss. Joannem et Paulum}}\jeune{Jeûne et abstinence}\dev{Treize vendredis de S. François de Paule — 6/13}{\pg{72}}}
\jourcal{Samedi}{13}{0}{violet}{\offpetit{Samedi après les Cendres}}{fér. maj.}{\liturgie{Intr. \textit{Audívit Dóminus} \enspace·\enspace \textit{Statio ad S. Tryphonem}}\jeune{Jeûne et abstinence}}
\notes
\semainegauche{du 14 au 20 février 2027}{1\textsuperscript{er} dimanche de Carême}
\jourcal{Dimanche}{14}{1}{violet}{\offgrand{1\textsuperscript{er} dimanche de Carême}}{dim. 1\textsuperscript{re} cl.}{\comm{S. Valentin, Prêtre et Martyr}\liturgie{Intr. \textit{Invocábit me} \enspace·\enspace \textit{Statio ad S. Joannem in Laterano}}\dev{Sept dimanches de S. Joseph — 3/7}{\pg{64}}}
\jourcal{Lundi}{15}{0}{violet}{\offpetit{Férie de Carême}}{fér. maj.}{\comm{Ss. Faustin et Jovite, Martyrs}\liturgie{Intr. \textit{Sicut óculi servórum in mánibus…} \enspace·\enspace \textit{Statio ad S. Petrum ad Vincula}}\jeune{Jeûne}}
\jourcal{Mardi}{16}{0}{violet}{\offpetit{Férie de Carême}}{fér. maj.}{\liturgie{Intr. \textit{Dómine} \enspace·\enspace \textit{Statio ad S. Anastasiam}}\jeune{Jeûne}}
\jourcal{Mercredi}{17}{0}{violet}{\offpetit{Mercredi des Quatre-Temps de Carême}}{fér. maj.}{\liturgie{Intr. \textit{Reminíscere miseratiónum tuárum} \enspace·\enspace \textit{Statio ad S. Mariam majorem}}\jeune{Jeûne et abstinence}}
\semainedroite{du 14 au 20 février 2027}
\jourcal{Jeudi}{18}{0}{violet}{\offpetit{Férie de Carême}}{fér. maj.}{\comm{S. Siméon, Évêque et Martyr}\liturgie{Intr. \textit{Conféssio et pulchritúdo in conspéctu…} \enspace·\enspace \textit{Statio ad S. Laurentium in Panisperna}}\jeune{Jeûne}}
\jourcal{Vendredi}{19}{0}{violet}{\offpetit{Vendredi des Quatre-Temps de Carême}}{fér. maj.}{\liturgie{Intr. \textit{De necessitátibus meis éripe me} \enspace·\enspace \textit{Statio ad Ss. duodecim Apostolos}}\jeune{Jeûne et abstinence}\dev{Treize vendredis de S. François de Paule — 7/13}{\pg{72}}}
\jourcal{Samedi}{20}{0}{violet}{\offpetit{Samedi des Quatre-Temps de Carême}}{fér. maj.}{\liturgie{Intr. \textit{Intret orátio mea in conspéctu…} \enspace·\enspace \textit{Statio ad S. Petrum}}\jeune{Jeûne et abstinence}}
\notes
\semainegauche{du 21 au 27 février 2027}{2\textsuperscript{e} dimanche de Carême}
\jourcal{Dimanche}{21}{1}{violet}{\offgrand{2\textsuperscript{e} dimanche de Carême}}{dim. 1\textsuperscript{re} cl.}{\liturgie{Intr. \textit{Reminíscere miseratiónum tuarum} \enspace·\enspace \textit{Statio ad S. Mariam in Domnica}}\dev{Sept dimanches de S. Joseph — 4/7}{\pg{64}}}
\jourcal{Lundi}{22}{0}{blanc}{\offgrand{Chaire de S. Pierre, Apôtre}}{dbl maj.}{\comm{Férie de Carême}\liturgie{Intr. \textit{Státuit ei Dóminus testaméntum pacis} \enspace·\enspace \textit{Statio ad S. Clementem}}\jeune{Jeûne}}
\jourcal{Mardi}{23}{0}{blanc}{\offpetit{S. Pierre Damien, Évêque, Confesseur et Docteur de l’Église}}{dbl}{\comm{Férie de Carême ; Vigile de S. Matthias, Apôtre}\liturgie{Intr. \textit{In médio Ecclésiæ apéruit os…} \enspace·\enspace \textit{Statio ad S. Balbinam}}\jeune{Jeûne}}
\jourfete{Mercredi}{24}{1}{rouge}{\offgrand{S. Matthias, Apôtre}}{2\textsuperscript{e} cl.}{\comm{Férie de Carême}\liturgie{Intr. \textit{Mihi autem nimis honoráti sunt…} \enspace·\enspace \textit{Statio ad S. Caeciliam}}\noticesaint{En Judée, l’anniversaire de saint Matthias Apôtre. Après l’Ascension du Seigneur, il fut choisi au sort par les Apôtres pour occuper la place du traître Judas. […]}\jeune{Jeûne}}
\semainedroite{du 21 au 27 février 2027}
\jourcal{Jeudi}{25}{0}{violet}{\offpetit{Férie de Carême}}{fér. maj.}{\liturgie{Intr. \textit{Deus} \enspace·\enspace \textit{Statio ad S. Mariam trans Tiberim}}\jeune{Jeûne}}
\jourcal{Vendredi}{26}{0}{violet}{\offpetit{Férie de Carême}}{fér. maj.}{\liturgie{Intr. \textit{Ego autem cum justítia apparébo…} \enspace·\enspace \textit{Statio ad S. Vitalem}}\jeune{Jeûne et abstinence}\dev{Treize vendredis de S. François de Paule — 8/13}{\pg{72}}}
\jourcal{Samedi}{27}{0}{blanc}{\offpetit{S. Gabriel de l’Addolorata (de N.-D. des Douleurs), Confesseur}}{dbl}{\comm{Férie de Carême}\liturgie{Intr. \textit{Oculus Dei respéxit illum in…} \enspace·\enspace \textit{Statio ad Ss. Marcellinum et Petrum}}\jeune{Jeûne et abstinence}}
\notes
\moisgauche{Mars}{2027}{\devmois{Mois de saint Joseph}{466}{0}}{Le juste fleurira comme le palmier et il se multipliera comme le cèdre du Liban, planté dans la maison du Seigneur, dans les parvis de la maison de notre Dieu.}{Ps 91:13-14}{S. Joseph}
\moisdroite{Mars 2027}{\rubmois{Fêtes principales}\lignemois{19}{S. Joseph, Époux de la Bienheureuse Vierge Marie, Confesseur}\lignemois{25}{Jeudi saint}\lignemois{26}{Vendredi saint}\lignemois{27}{Samedi saint}\lignemois{28}{Dimanche de Pâques}\lignemois{29}{Lundi de Pâques}\lignemois{30}{Mardi de Pâques}\lignemois{31}{Mercredi de Pâques}\rubmois{Dévotions et neuvaines qui commencent}\lignemois{3}{Montfort, douze jours : se vider de l’esprit du monde\pg{11}}\lignemois{3}{Préparation à la consécration de Montfort\pg{11}}\lignemois{10}{Neuvaine à S. Joseph\pg{34}}\lignemois{13}{Préparation à la consécration de Montfort\pg{11}}\lignemois{15}{Montfort, 1\textsuperscript{re} semaine : connaissance de soi\pg{11}}\lignemois{15}{Neuvaine à S. Gabriel\pg{19}}\lignemois{16}{Treize mardis de S. Antoine\pg{68}}\lignemois{22}{Montfort, 2\textsuperscript{e} semaine : connaissance de la sainte Vierge\pg{11}}\lignemois{29}{Montfort, 3\textsuperscript{e} semaine : connaissance de Jésus-Christ\pg{11}}\rubmois{Pénitence}\lignemoislong{Jeûne et abstinence}{5, 6, 12, 13, 19, 20, 26, 27}\lignemoislong{Jeûne}{1, 2, 3, 4, 8, 9, 10, 11, 15, 16, 17, 18, 22, 23, 24, 25}}
\semainegauche{du 28 février au 6 mars 2027}{3\textsuperscript{e} dimanche de Carême}
\jourcal{Dimanche}{28}{1}{violet}{\offgrand{3\textsuperscript{e} dimanche de Carême}}{dim. 1\textsuperscript{re} cl.}{\liturgie{Intr. \textit{Oculi mei semper ad Dóminum} \enspace·\enspace \textit{Statio ad S. Laurentium extra muros}}\dev{Sept dimanches de S. Joseph — 5/7}{\pg{64}}}
\jourcal{Lundi}{1}{0}{violet}{\offpetit{Férie de Carême}}{fér. maj.}{\liturgie{Intr. \textit{In Deo laudábo verbum} \enspace·\enspace \textit{Statio ad S. Marcum}}\jeune{Jeûne}}
\jourcal{Mardi}{2}{0}{violet}{\offpetit{Férie de Carême}}{fér. maj.}{\liturgie{Intr. \textit{Ego clamávi} \enspace·\enspace \textit{Statio ad S. Pudentianam}}\jeune{Jeûne}}
\jourcal{Mercredi}{3}{0}{violet}{\offpetit{Férie de Carême}}{fér. maj.}{\liturgie{Intr. \textit{Ego autem in Dómino sperábo} \enspace·\enspace \textit{Statio ad S. Xystum}}\jeune{Jeûne}\dev{Premier mercredi : S. Joseph}{\pg{68}}\dev{Montfort, douze jours : se vider de l’esprit du monde — 1/12}{\pg{11}}}
\semainedroite{du 28 février au 6 mars 2027}
\jourcal{Jeudi}{4}{0}{blanc}{\offpetit{S. Casimir, Confesseur}}{semi-dbl}{\comm{Férie de Carême ; S. Lucien I\textsuperscript{er}, Pape et Martyr}\liturgie{Intr. \textit{Os justi meditábitur sapiéntiam} \enspace·\enspace \textit{Statio ad Ss. Cosmam et Damianum}}\jeune{Jeûne}\dev{Heure sainte (veille du premier vendredi)}{\pg{69}}\dev{Montfort, douze jours : se vider de l’esprit du monde — 2/12}{\pg{11}}}
\jourcal{Vendredi}{5}{0}{violet}{\offpetit{Férie de Carême}}{fér. maj.}{\liturgie{Intr. \textit{Fac mecum} \enspace·\enspace \textit{Statio ad S. Laurentium in Lucina}}\jeune{Jeûne et abstinence}\dev{Treize vendredis de S. François de Paule — 9/13}{\pg{72}}\dev{Premier vendredi du mois — 4/9}{\pg{69}}\dev{Montfort, douze jours : se vider de l’esprit du monde — 3/12}{\pg{11}}}
\jourcal{Samedi}{6}{0}{rouge}{\offpetit{Stes Perpétue et Félicité, Martyres}}{dbl}{\comm{Férie de Carême}\liturgie{Intr. \textit{Me exspectavérunt peccatóres} \enspace·\enspace \textit{Statio ad S. Susannam}}\jeune{Jeûne et abstinence}\dev{Premier samedi du mois — 4/5}{\pg{73}}\dev{Montfort, douze jours : se vider de l’esprit du monde — 4/12}{\pg{11}}}
\notes
\semainegauche{du 7 au 13 mars 2027}{4\textsuperscript{e} dimanche de Carême (Laetare)}
\jourcal{Dimanche}{7}{1}{rose}{\offgrand{4\textsuperscript{e} dimanche de Carême (Laetare)}}{dim. 1\textsuperscript{re} cl.}{\comm{S. Thomas d’Aquin, Confesseur et Docteur de l’Église}\liturgie{Intr. \textit{Lætáre} \enspace·\enspace \textit{Statio ad S. Crucem in Jerusalem}}\dev{Sept dimanches de S. Joseph — 6/7}{\pg{64}}\dev{Montfort, douze jours : se vider de l’esprit du monde — 5/12}{\pg{11}}}
\jourcal{Lundi}{8}{0}{blanc}{\offpetit{S. Jean de Dieu, Confesseur}}{dbl}{\comm{Férie de Carême}\liturgie{Intr. \textit{Os justi meditábitur sapiéntiam} \enspace·\enspace \textit{Statio ad Ss. Quatuor Coronatos}}\jeune{Jeûne}\dev{Montfort, douze jours : se vider de l’esprit du monde — 6/12}{\pg{11}}}
\jourcal{Mardi}{9}{0}{blanc}{\offpetit{Ste Françoise Romaine, Veuve}}{dbl}{\comm{Férie de Carême}\liturgie{Intr. \textit{Cognóvi} \enspace·\enspace \textit{Statio ad S. Laurentium in Damaso}}\jeune{Jeûne}\dev{Montfort, douze jours : se vider de l’esprit du monde — 7/12}{\pg{11}}}
\jourcal{Mercredi}{10}{0}{rouge}{\offpetit{Les Quarante Martyrs}}{semi-dbl}{\comm{Férie de Carême}\liturgie{Intr. \textit{Clamavérunt justi} \enspace·\enspace \textit{Statio ad S. Paulum}}\jeune{Jeûne}\dev{Montfort, douze jours : se vider de l’esprit du monde — 8/12}{\pg{11}}\nov{Neuvaine à S. Joseph — 1\textsuperscript{er} jour}{\pg{34}}}
\semainedroite{du 7 au 13 mars 2027}
\jourcal{Jeudi}{11}{0}{violet}{\offpetit{Férie de Carême}}{fér. maj.}{\liturgie{Intr. \textit{Lætétur cor quæréntium Dóminum} \enspace·\enspace \textit{Statio ad Ss. Silvestrum et Martinum}}\jeune{Jeûne}\dev{Montfort, douze jours : se vider de l’esprit du monde — 9/12}{\pg{11}}\nov{Neuvaine à S. Joseph — 2\textsuperscript{e} jour}{\pg{34}}}
\jourcal{Vendredi}{12}{0}{blanc}{\offpetit{S. Grégoire le Grand, Pape, Confesseur et Docteur de l’Église}}{dbl}{\comm{Férie de Carême}\liturgie{Intr. \textit{Státuit ei Dóminus testaméntum pacis} \enspace·\enspace \textit{Statio ad S. Eusebium}}\jeune{Jeûne et abstinence}\dev{Treize vendredis de S. François de Paule — 10/13}{\pg{72}}\dev{Montfort, douze jours : se vider de l’esprit du monde — 10/12}{\pg{11}}\nov{Neuvaine à S. Joseph — 3\textsuperscript{e} jour}{\pg{34}}}
\jourcal{Samedi}{13}{0}{violet}{\offpetit{Férie de Carême}}{fér. maj.}{\liturgie{Intr. \textit{Sitiéntes} \enspace·\enspace \textit{Statio ad S. Nicolaum in Carcere}}\jeune{Jeûne et abstinence}\dev{Montfort, douze jours : se vider de l’esprit du monde — 11/12}{\pg{11}}\nov{Neuvaine à S. Joseph — 4\textsuperscript{e} jour}{\pg{34}}}
\notes
\semainegauche{du 14 au 20 mars 2027}{1\textsuperscript{er} dimanche de la Passion}
\jourcal{Dimanche}{14}{1}{violet}{\offgrand{1\textsuperscript{er} dimanche de la Passion}}{dim. 1\textsuperscript{re} cl.}{\liturgie{Intr. \textit{Júdica me} \enspace·\enspace \textit{Statio ad S. Petrum}}\dev{Sept dimanches de S. Joseph — 7/7}{\pg{64}}\dev{Montfort, douze jours : se vider de l’esprit du monde — 12/12}{\pg{11}}\nov{Neuvaine à S. Joseph — 5\textsuperscript{e} jour}{\pg{34}}}
\jourcal{Lundi}{15}{0}{violet}{\offpetit{Lundi de la Passion}}{fér. maj.}{\liturgie{Intr. \textit{Miserére mihi} \enspace·\enspace \textit{Statio ad S. Chrysogonum}}\jeune{Jeûne}\dev{Montfort, 1\textsuperscript{re} semaine : connaissance de soi — 1/7}{\pg{11}}\nov{Neuvaine à S. Joseph — 6\textsuperscript{e} jour}{\pg{34}}\nov{Neuvaine à S. Gabriel — 1\textsuperscript{er} jour}{\pg{19}}}
\jourcal{Mardi}{16}{0}{violet}{\offpetit{Mardi de la Passion}}{fér. maj.}{\liturgie{Intr. \textit{Exspécta Dóminum} \enspace·\enspace \textit{Statio ad S. Cyriacum}}\jeune{Jeûne}\dev{Treize mardis de S. Antoine — 1/13}{\pg{68}}\dev{Montfort, 1\textsuperscript{re} semaine : connaissance de soi — 2/7}{\pg{11}}\nov{Neuvaine à S. Joseph — 7\textsuperscript{e} jour}{\pg{34}}\nov{Neuvaine à S. Gabriel — 2\textsuperscript{e} jour}{\pg{19}}}
\jourcal{Mercredi}{17}{0}{blanc}{\offpetit{S. Patrick, Évêque et Confesseur}}{dbl}{\comm{Mercredi de la Passion}\liturgie{Intr. \textit{Státuit ei Dóminus testaméntum pacis} \enspace·\enspace \textit{Statio ad S. Marcellum}}\jeune{Jeûne}\dev{Montfort, 1\textsuperscript{re} semaine : connaissance de soi — 3/7}{\pg{11}}\nov{Neuvaine à S. Joseph — 8\textsuperscript{e} jour}{\pg{34}}\nov{Neuvaine à S. Gabriel — 3\textsuperscript{e} jour}{\pg{19}}}
\semainedroite{du 14 au 20 mars 2027}
\jourcal{Jeudi}{18}{0}{blanc}{\offpetit{S. Cyrille, Évêque de Jérusalem, Confesseur et Docteur de l’Église}}{dbl}{\comm{Jeudi de la Passion}\liturgie{Intr. \textit{In médio Ecclésiæ apéruit os…} \enspace·\enspace \textit{Statio ad S. Apollinarem}}\jeune{Jeûne}\dev{Montfort, 1\textsuperscript{re} semaine : connaissance de soi — 4/7}{\pg{11}}\nov{Neuvaine à S. Joseph — 9\textsuperscript{e} jour}{\pg{34}}\nov{Neuvaine à S. Gabriel — 4\textsuperscript{e} jour}{\pg{19}}}
\jourfete{Vendredi}{19}{1}{blanc}{\offgrand{S. Joseph, Époux de la Bienheureuse Vierge Marie, Confesseur}}{1\textsuperscript{re} cl.}{\comm{Vendredi de la Passion}\liturgie{Intr. \textit{Justus ut palma florébit} \enspace·\enspace \textit{Statio ad S. Stephanum in Cælio monte}}\noticesaint{En Judée, l’anniversaire de saint Joseph, époux de la Bienheureuse Vierge Marie, Confesseur et Patron de l’Église universelle. […]}\jeune{Jeûne et abstinence}\dev{Treize vendredis de S. François de Paule — 11/13}{\pg{72}}\dev{Montfort, 1\textsuperscript{re} semaine : connaissance de soi — 5/7}{\pg{11}}\dev{Douleurs et Allégresses de S. Joseph ; hymne Te, Joseph, celebrent (p. 60)}{\pg{34}}\nov{Neuvaine à S. Gabriel — 5\textsuperscript{e} jour}{\pg{19}}}
\jourcal{Samedi}{20}{0}{violet}{\offpetit{Samedi de la Passion}}{fér. maj.}{\liturgie{Intr. \textit{Miserére mihi} \enspace·\enspace \textit{Statio ad S. Joannem ante Portam Latinam}}\jeune{Jeûne et abstinence}\dev{Montfort, 1\textsuperscript{re} semaine : connaissance de soi — 6/7}{\pg{11}}\nov{Neuvaine à S. Gabriel — 6\textsuperscript{e} jour}{\pg{19}}}
\notes
\semainegauche{du 21 au 27 mars 2027}{Dimanche des Rameaux}
\jourcal{Dimanche}{21}{1}{violet}{\offgrand{Dimanche des Rameaux}}{dim. 1\textsuperscript{re} cl.}{\comm{S. Benoît, Abbé}\liturgie{Intr. \textit{Dómine} \enspace·\enspace \textit{Statio ad S. Joannem in Laterano}}\dev{Montfort, 1\textsuperscript{re} semaine : connaissance de soi — 7/7}{\pg{11}}\nov{Neuvaine à S. Gabriel — 7\textsuperscript{e} jour}{\pg{19}}}
\jourcal{Lundi}{22}{0}{violet}{\offpetit{Lundi saint}}{fér. priv.}{\liturgie{Intr. \textit{Júdica} \enspace·\enspace \textit{Statio ad S. Praxedem}}\jeune{Jeûne}\dev{Montfort, 2\textsuperscript{e} semaine : connaissance de la sainte Vierge — 1/7}{\pg{11}}\nov{Neuvaine à S. Gabriel — 8\textsuperscript{e} jour}{\pg{19}}}
\jourcal{Mardi}{23}{0}{violet}{\offpetit{Mardi saint}}{fér. priv.}{\liturgie{Intr. \textit{Nos autem gloriári opórtet in…} \enspace·\enspace \textit{Statio ad S. Priscam}}\jeune{Jeûne}\dev{Treize mardis de S. Antoine — 2/13}{\pg{68}}\dev{Montfort, 2\textsuperscript{e} semaine : connaissance de la sainte Vierge — 2/7}{\pg{11}}\nov{Neuvaine à S. Gabriel — 9\textsuperscript{e} jour}{\pg{19}}}
\jourcal{Mercredi}{24}{0}{violet}{\offpetit{Mercredi saint}}{fér. priv.}{\liturgie{Intr. \textit{In nómine Jesu omne genu…} \enspace·\enspace \textit{Statio ad S. Mariam majorem}}\jeune{Jeûne}\dev{Montfort, 2\textsuperscript{e} semaine : connaissance de la sainte Vierge — 3/7}{\pg{11}}\dev{Prière à S. Gabriel}{\pg{19}}}
\semainedroite{du 21 au 27 mars 2027}
\jourfete{Jeudi}{25}{1}{blanc}{\offgrand{Jeudi saint}}{1\textsuperscript{re} cl.}{\liturgie{Intr. \textit{Nos autem gloriári opórtet in…} \enspace·\enspace \textit{Statio ad S. Joannem in Laterano}}\jeune{Jeûne}\dev{Montfort, 2\textsuperscript{e} semaine : connaissance de la sainte Vierge — 4/7}{\pg{11}}}
\jourfete{Vendredi}{26}{1}{noir}{\offgrand{Vendredi saint}}{1\textsuperscript{re} cl.}{\liturgie{Intr. \textit{Dómine} \enspace·\enspace \textit{Statio ad S. Crucem in Jerusalem}}\jeune{Jeûne et abstinence}\dev{Treize vendredis de S. François de Paule — 12/13}{\pg{72}}\dev{Montfort, 2\textsuperscript{e} semaine : connaissance de la sainte Vierge — 5/7}{\pg{11}}}
\jourfete{Samedi}{27}{1}{violet}{\offgrand{Samedi saint}}{1\textsuperscript{re} cl.}{\liturgie{Intr. \textit{Dómine} \enspace·\enspace \textit{Statio ad S. Joannem in Laterano}}\jeune{Jeûne et abstinence jusqu’à midi}\dev{Montfort, 2\textsuperscript{e} semaine : connaissance de la sainte Vierge — 6/7}{\pg{11}}}
\notes
\moisgauche{Avril}{2027}{}{Je suis ressuscité, et je suis de nouveau avec vous, alléluia ; vous avez placé sur moi votre main, alléluia ; votre science est admirable, alléluia, alléluia.}{Ps 138:18; 138:5-6}{Pâques}
\moisdroite{Avril 2027}{\rubmois{Fêtes principales}\lignemois{1}{Jeudi de Pâques}\lignemois{2}{Vendredi de Pâques}\lignemois{3}{Samedi de Pâques}\lignemois{5}{Annonciation de la Bienheureuse Vierge Marie}\lignemois{14}{Patronage de S. Joseph, époux de la Bse Vierge Marie}\lignemois{25}{S. Marc, Évangéliste}\rubmois{Dévotions et neuvaines qui commencent}\lignemois{11}{Six dimanches de S. Louis de Gonzague\pg{63}}\lignemois{17}{Quinze samedis du Rosaire\pg{74}}\lignemois{19}{Neuvaine à S. Louis-Marie Grignion de Montfort (28 avr.)}\rubmois{Pénitence}\lignemoislong{Abstinence}{2, 9, 16, 23, 30}}
\semainegauche{du 28 mars au 3 avril 2027}{Dimanche de Pâques}
\jourcal{Dimanche}{28}{1}{blanc}{\offgrand{Dimanche de Pâques}}{1\textsuperscript{re} cl.}{\liturgie{Intr. \textit{Resurréxi} \enspace·\enspace \textit{Statio ad S. Mariam majorem}}\dev{Montfort, 2\textsuperscript{e} semaine : connaissance de la sainte Vierge — 7/7}{\pg{11}}}
\jourfete{Lundi}{29}{1}{blanc}{\offgrand{Lundi de Pâques}}{1\textsuperscript{re} cl.}{\liturgie{Intr. \textit{Introdúxit vos Dóminus in terram…} \enspace·\enspace \textit{Statio ad S. Petrum}}\dev{Montfort, 3\textsuperscript{e} semaine : connaissance de Jésus-Christ — 1/7}{\pg{11}}}
\jourfete{Mardi}{30}{1}{blanc}{\offgrand{Mardi de Pâques}}{1\textsuperscript{re} cl.}{\liturgie{Intr. \textit{Aqua sapiéntiæ potávit eos} \enspace·\enspace \textit{Statio ad S. Paulum}}\dev{Treize mardis de S. Antoine — 3/13}{\pg{68}}\dev{Montfort, 3\textsuperscript{e} semaine : connaissance de Jésus-Christ — 2/7}{\pg{11}}}
\jourfete{Mercredi}{31}{1}{blanc}{\offgrand{Mercredi de Pâques}}{1\textsuperscript{re} cl.}{\liturgie{Intr. \textit{Veníte} \enspace·\enspace \textit{Statio ad S. Laurentium extra muros}}\dev{Montfort, 3\textsuperscript{e} semaine : connaissance de Jésus-Christ — 3/7}{\pg{11}}}
\semainedroite{du 28 mars au 3 avril 2027}
\jourfete{Jeudi}{1}{1}{blanc}{\offgrand{Jeudi de Pâques}}{1\textsuperscript{re} cl.}{\liturgie{Intr. \textit{Victrícem manum tuam} \enspace·\enspace \textit{Statio ad Ss. duodecim Apostolos}}\dev{Heure sainte (veille du premier vendredi)}{\pg{69}}\dev{Montfort, 3\textsuperscript{e} semaine : connaissance de Jésus-Christ — 4/7}{\pg{11}}}
\jourfete{Vendredi}{2}{1}{blanc}{\offgrand{Vendredi de Pâques}}{1\textsuperscript{re} cl.}{\liturgie{Intr. \textit{Edúxit eos Dóminus in spe} \enspace·\enspace \textit{Statio ad S. Mariam ad Martyres}}\jeune{Abstinence}\dev{Treize vendredis de S. François de Paule — 13/13}{\pg{72}}\dev{Premier vendredi du mois — 5/9}{\pg{69}}\dev{Montfort, 3\textsuperscript{e} semaine : connaissance de Jésus-Christ — 5/7}{\pg{11}}}
\jourfete{Samedi}{3}{1}{blanc}{\offgrand{Samedi de Pâques}}{1\textsuperscript{re} cl.}{\liturgie{Intr. \textit{Edúxit Dóminus pópulum suum in…} \enspace·\enspace \textit{Statio ad S. Joannem in Laterano}}\dev{Premier samedi du mois — 5/5}{\pg{73}}\dev{Montfort, 3\textsuperscript{e} semaine : connaissance de Jésus-Christ — 6/7}{\pg{11}}}
\notes
\semainegauche{du 4 au 10 avril 2027}{Dimanche in Albis (Quasimodo)}
\jourcal{Dimanche}{4}{1}{blanc}{\offgrand{Dimanche in Albis (Quasimodo)}}{dim. 1\textsuperscript{re} cl.}{\comm{S. Isidore, Évêque, Confesseur et Docteur}\liturgie{Intr. \textit{Quasi modo géniti infántes} \enspace·\enspace \textit{Statio ad S. Pancratium}}\dev{Montfort, 3\textsuperscript{e} semaine : connaissance de Jésus-Christ — 7/7}{\pg{11}}}
\jourfete{Lundi}{5}{1}{blanc}{\offgrand{Annonciation de la Bienheureuse Vierge Marie}}{1\textsuperscript{re} cl.}{\comm{S. Vincent Ferrier, Confesseur}\liturgie{Intr. \textit{Vultum tuum deprecabúntur omnes dívites…}}\noticesaint{L’Annonciation de la bienheureuse Vierge Marie, Mère de Dieu.}\dev{Consécration (renouvellement) à Jésus par Marie — communion de ce jour}{\pg{11}}}
\jourcal{Mardi}{6}{0}{blanc}{\offpetit{Férie du temps pascal}}{férie}{\liturgie{Intr. \textit{Quasi modo géniti infántes}}\dev{Treize mardis de S. Antoine — 4/13}{\pg{68}}}
\jourcal{Mercredi}{7}{0}{blanc}{\offpetit{Férie du temps pascal}}{férie}{\liturgie{Intr. \textit{Quasi modo géniti infántes}}\dev{Premier mercredi : S. Joseph}{\pg{68}}}
\semainedroite{du 4 au 10 avril 2027}
\jourcal{Jeudi}{8}{0}{blanc}{\offpetit{Férie du temps pascal}}{férie}{\liturgie{Intr. \textit{Quasi modo géniti infántes}}}
\jourcal{Vendredi}{9}{0}{blanc}{\offpetit{Férie du temps pascal}}{férie}{\liturgie{Intr. \textit{Quasi modo géniti infántes}}\jeune{Abstinence}}
\jourcal{Samedi}{10}{0}{blanc}{\offpetit{Office de la Sainte Vierge le samedi}}{simple}{\liturgie{Intr. \textit{Salve, sancta Parens}}}
\notes
\semainegauche{du 11 au 17 avril 2027}{2\textsuperscript{e} dimanche après Pâques}
\jourcal{Dimanche}{11}{1}{blanc}{\offgrand{2\textsuperscript{e} dimanche après Pâques}}{dim.}{\comm{S. Léon le Grand, Pape, Confesseur et Docteur}\liturgie{Intr. \textit{Misericórdia Dómini plena est terra}}\dev{Six dimanches de S. Louis de Gonzague — 1/6}{\pg{63}}}
\jourcal{Lundi}{12}{0}{blanc}{\offpetit{Férie du temps pascal}}{férie}{\liturgie{Intr. \textit{Misericórdia Dómini plena est terra}}}
\jourcal{Mardi}{13}{0}{rouge}{\offpetit{S. Herménégilde, Martyr}}{semi-dbl}{\liturgie{Intr. \textit{In virtúte tua}}\dev{Treize mardis de S. Antoine — 5/13}{\pg{68}}}
\jourfete{Mercredi}{14}{1}{blanc}{\offgrand{Patronage de S. Joseph, époux de la Bse Vierge Marie}}{1\textsuperscript{re} cl.}{\comm{S. Justin, Martyr ; Ss. Tiburce, Valérien et Maxime, Martyrs}\liturgie{Intr. \textit{Adjútor et protéctor noster est…}}}
\semainedroite{du 11 au 17 avril 2027}
\jourcal{Jeudi}{15}{0}{blanc}{\offpetit{Deuxième jour dans l’octave de S. Joseph}}{semi-dbl}{\liturgie{Intr. \textit{Adjútor et protéctor noster est…}}}
\jourcal{Vendredi}{16}{0}{blanc}{\offpetit{Troisième jour dans l’octave de S. Joseph}}{semi-dbl}{\liturgie{Intr. \textit{Adjútor et protéctor noster est…}}\jeune{Abstinence}}
\jourcal{Samedi}{17}{0}{blanc}{\offpetit{Quatrième jour dans l’octave de S. Joseph}}{semi-dbl}{\comm{S. Anicet, Pape et Martyr}\liturgie{Intr. \textit{Adjútor et protéctor noster est…}}\dev{Quinze samedis du Rosaire — 1/15}{\pg{74}}}
\notes
\semainegauche{du 18 au 24 avril 2027}{3\textsuperscript{e} dimanche après Pâques}
\jourcal{Dimanche}{18}{1}{blanc}{\offgrand{3\textsuperscript{e} dimanche après Pâques}}{dim.}{\comm{Octave de S. Joseph}\liturgie{Intr. \textit{Jubiláte Deo}}\dev{Six dimanches de S. Louis de Gonzague — 2/6}{\pg{63}}}
\jourcal{Lundi}{19}{0}{blanc}{\offpetit{Sixième jour dans l’octave de S. Joseph}}{semi-dbl}{\liturgie{Intr. \textit{Adjútor et protéctor noster est…}}\nov{Neuvaine à S. Louis-Marie Grignion de Montfort (28 avr.) — 1\textsuperscript{er} jour}{}}
\jourcal{Mardi}{20}{0}{blanc}{\offpetit{Septième jour dans l’octave de S. Joseph}}{semi-dbl}{\liturgie{Intr. \textit{Adjútor et protéctor noster est…}}\dev{Treize mardis de S. Antoine — 6/13}{\pg{68}}\nov{Neuvaine à S. Louis-Marie Grignion de Montfort (28 avr.) — 2\textsuperscript{e} jour}{}}
\jourcal{Mercredi}{21}{0}{blanc}{\offgrand{Octave de S. Joseph}}{dbl maj.}{\comm{S. Anselme, Évêque, Confesseur et Docteur}\liturgie{Intr. \textit{Adjútor et protéctor noster est…}}\nov{Neuvaine à S. Louis-Marie Grignion de Montfort (28 avr.) — 3\textsuperscript{e} jour}{}}
\semainedroite{du 18 au 24 avril 2027}
\jourcal{Jeudi}{22}{0}{rouge}{\offpetit{Ss. Soter et Caius, Papes et Martyrs}}{semi-dbl}{\liturgie{Intr. \textit{Intret in conspéctu tuo}}\nov{Neuvaine à S. Louis-Marie Grignion de Montfort (28 avr.) — 4\textsuperscript{e} jour}{}}
\jourcal{Vendredi}{23}{0}{rouge}{\offpetit{S. Georges, Martyr}}{semi-dbl}{\liturgie{Intr. \textit{In virtúte tua}}\jeune{Abstinence}\nov{Neuvaine à S. Louis-Marie Grignion de Montfort (28 avr.) — 5\textsuperscript{e} jour}{}}
\jourcal{Samedi}{24}{0}{rouge}{\offpetit{S. Fidèle de Sigmaringen, Martyr}}{dbl}{\liturgie{Intr. \textit{In virtúte tua}}\dev{Quinze samedis du Rosaire — 2/15}{\pg{74}}\nov{Neuvaine à S. Louis-Marie Grignion de Montfort (28 avr.) — 6\textsuperscript{e} jour}{}}
\notes
\moisgauche{Mai}{2027}{\devmois{Mois de Marie}{325}{1}}{L'Esprit du Seigneur remplit la terre, alléluia ; et lui qui tient unies toutes choses perce le secret de chaque parole, alléluia, alléluia, alléluia.}{Sap 1:7}{Pentecôte}
\moisdroite{Mai 2027}{\rubmois{Fêtes principales}\lignemois{1}{Ss. Philippe et Jacques, Apôtres}\lignemois{3}{Invention de la Sainte Croix}\lignemois{6}{Ascension de Notre-Seigneur}\lignemois{8}{L’Apparition de S. Michel Archange}\lignemois{16}{Dimanche de la Pentecôte}\lignemois{17}{Lundi de Pentecôte}\lignemois{18}{Mardi de Pentecôte}\lignemois{19}{Mercredi des Quatre-Temps de Pentecôte}\lignemois{20}{Jeudi de Pentecôte}\lignemois{21}{Vendredi des Quatre-Temps de Pentecôte}\lignemois{22}{Samedi des Quatre-Temps de Pentecôte}\lignemois{23}{Très Sainte Trinité}\lignemois{27}{Fête-Dieu (Saint-Sacrement)}\lignemois{28}{Dans l’octave du Saint-Sacrement}\lignemois{29}{Dans l’octave du Saint-Sacrement}\lignemois{31}{Dans l’octave du Saint-Sacrement}\rubmois{Dévotions et neuvaines qui commencent}\lignemois{7}{Neuvaine au Saint-Esprit, avant la Pentecôte\pg{56}}\lignemois{23}{Dix dimanches de S. Ignace\pg{66}}\lignemois{26}{Neuvaine au Sacré-Cœur}\rubmois{Pénitence}\lignemoislong{Jeûne et abstinence}{15, 19, 21, 22}\lignemoislong{Abstinence}{7, 14, 28}}
\semainegauche{du 25 avril au 1 mai 2027}{4\textsuperscript{e} dimanche après Pâques}
\jourfete{Dimanche}{25}{1}{rouge}{\offgrand{S. Marc, Évangéliste}}{2\textsuperscript{e} cl.}{\comm{4\textsuperscript{e} dimanche après Pâques}\liturgie{Intr. \textit{Protexísti me} \enspace·\enspace \textit{Litaniæ majores : Statio ad S. Petrum}}\noticesaint{À Alexandrie, l’anniversaire du Bienheureux Évangéliste Marc. […]}\dev{Six dimanches de S. Louis de Gonzague — 3/6}{\pg{63}}\nov{Neuvaine à S. Louis-Marie Grignion de Montfort (28 avr.) — 7\textsuperscript{e} jour}{}}
\jourcal{Lundi}{26}{0}{rouge}{\offpetit{Ss. Clet et Marcellin, Papes et Martyrs}}{semi-dbl}{\liturgie{Intr. \textit{Intret in conspéctu tuo}}\nov{Neuvaine à S. Louis-Marie Grignion de Montfort (28 avr.) — 8\textsuperscript{e} jour}{}}
\jourcal{Mardi}{27}{0}{blanc}{\offpetit{S. Pierre Canisius, Confesseur et Docteur de l’Église}}{dbl}{\liturgie{Intr. \textit{In médio Ecclésiæ apéruit os…}}\dev{Treize mardis de S. Antoine — 7/13}{\pg{68}}\nov{Neuvaine à S. Louis-Marie Grignion de Montfort (28 avr.) — 9\textsuperscript{e} jour}{}}
\jourcal{Mercredi}{28}{0}{rouge}{\offpetit{S. Paul de la Croix, Confesseur}}{dbl}{\comm{S. Vital, Martyr}\liturgie{Intr. \textit{Christo confíxus sum Cruci}}}
\semainedroite{du 25 avril au 1 mai 2027}
\jourcal{Jeudi}{29}{0}{rouge}{\offpetit{S. Pierre, Martyr}}{dbl}{\liturgie{Intr. \textit{Lætábitur justus in Dómino}}}
\jourcal{Vendredi}{30}{0}{blanc}{\offpetit{Ste Catherine de Sienne, Vierge}}{dbl}{\liturgie{Intr. \textit{Dilexísti justítiam}}\jeune{Abstinence}}
\jourfete{Samedi}{1}{1}{rouge}{\offgrand{Ss. Philippe et Jacques, Apôtres}}{2\textsuperscript{e} cl.}{\liturgie{Intr. \textit{Clamavérunt ad te}}\noticesaint{L’anniversaire des bienheureux Apôtres Philippe et Jacques. […]}\dev{Quinze samedis du Rosaire — 3/15}{\pg{74}}\dev{Premier samedi du mois}{\pg{73}}}
\notes
\semainegauche{du 2 au 8 mai 2027}{5\textsuperscript{e} dimanche après Pâques}
\jourcal{Dimanche}{2}{1}{blanc}{\offgrand{5\textsuperscript{e} dimanche après Pâques}}{dim.}{\comm{S. Athanase, Évêque, Confesseur et Docteur de l’Église}\liturgie{Intr. \textit{Vocem jucunditátis annuntiáte}}\dev{Six dimanches de S. Louis de Gonzague — 4/6}{\pg{63}}}
\jourfete{Lundi}{3}{1}{rouge}{\offgrand{Invention de la Sainte Croix}}{2\textsuperscript{e} cl.}{\comm{Lundi des Rogations ; Ss. Alexandre I, Pape, Évence et Théodule, Martyrs, et Juvénal, Évêque et Confesseur}\liturgie{Intr. \textit{Nos autem gloriári opórtet in…} \enspace·\enspace \textit{Statio ad S. Mariam majorem}}\noticesaint{À Jérusalem, l’Invention de la très sainte Croix du Seigneur, sous l’empereur Constantin.}}
\jourcal{Mardi}{4}{0}{blanc}{\offpetit{Ste Monique, Veuve}}{dbl}{\liturgie{Intr. \textit{Cognóvi} \enspace·\enspace \textit{Statio ad S. Joannem in Laterano}}\dev{Treize mardis de S. Antoine — 8/13}{\pg{68}}}
\jourcal{Mercredi}{5}{0}{blanc}{\offpetit{S. Pie V, Pape et Confesseur}}{dbl}{\comm{Mercredi des Rogations (vigile de l’Ascension)}\liturgie{Intr. \textit{Státuit ei Dóminus testaméntum pacis} \enspace·\enspace \textit{Statio ad S. Petrum}}\dev{Premier mercredi : S. Joseph}{\pg{68}}}
\semainedroite{du 2 au 8 mai 2027}
\jourfete{Jeudi}{6}{1}{blanc}{\offgrand{Ascension de Notre-Seigneur}}{1\textsuperscript{re} cl.}{\comm{S. Jean devant la Porte Latine}\liturgie{Intr. \textit{Viri Galilǽi} \enspace·\enspace \textit{Statio ad S. Petrum}}\dev{Heure sainte (veille du premier vendredi)}{\pg{69}}}
\jourcal{Vendredi}{7}{0}{rouge}{\offpetit{S. Stanislas, Évêque et Martyr}}{dbl}{\comm{Dans l’octave de l’Ascension}\liturgie{Intr. \textit{Sacerdótes Dei}}\jeune{Abstinence}\dev{Premier vendredi du mois — 6/9}{\pg{69}}\nov{Neuvaine au Saint-Esprit, avant la Pentecôte — 1\textsuperscript{er} jour}{\pg{56}}}
\jourcal{Samedi}{8}{0}{blanc}{\offgrand{L’Apparition de S. Michel Archange}}{dbl maj.}{\comm{Dans l’octave de l’Ascension}\liturgie{Intr. \textit{Benedícite Dóminum}}\noticesaint{Au mont Gargan, l’apparition de saint Michel Archange. Le Pape Pie XII l’a établi Patron et Protecteur des radiologues et de ceux qui soignent par le radium.}\dev{Quinze samedis du Rosaire — 4/15}{\pg{74}}\nov{Neuvaine au Saint-Esprit, avant la Pentecôte — 2\textsuperscript{e} jour}{\pg{56}}}
\notes
\semainegauche{du 9 au 15 mai 2027}{Dimanche dans l’octave de l’Ascension}
\jourcal{Dimanche}{9}{1}{blanc}{\offgrand{Dimanche dans l’octave de l’Ascension}}{dim.}{\comm{S. Grégoire de Nazianze, Évêque, Confesseur et Docteur de l’Église}\liturgie{Intr. \textit{Exáudi}}\dev{Six dimanches de S. Louis de Gonzague — 5/6}{\pg{63}}\nov{Neuvaine au Saint-Esprit, avant la Pentecôte — 3\textsuperscript{e} jour}{\pg{56}}}
\jourcal{Lundi}{10}{0}{blanc}{\offpetit{S. Antonin, Évêque et Confesseur}}{dbl}{\comm{Dans l’octave de l’Ascension ; Ss. Gordien et Épimaque, Martyrs}\liturgie{Intr. \textit{Státuit ei Dóminus testaméntum pacis}}\nov{Neuvaine au Saint-Esprit, avant la Pentecôte — 4\textsuperscript{e} jour}{\pg{56}}}
\jourcal{Mardi}{11}{0}{blanc}{\offpetit{Dans l’octave de l’Ascension}}{semi-dbl}{\liturgie{Intr. \textit{Viri Galilǽi}}\dev{Treize mardis de S. Antoine — 9/13}{\pg{68}}\nov{Neuvaine au Saint-Esprit, avant la Pentecôte — 5\textsuperscript{e} jour}{\pg{56}}}
\jourcal{Mercredi}{12}{0}{rouge}{\offpetit{Ss. Nérée, Achille, Domitille, Vierge, et Pancrace, Martyrs}}{semi-dbl}{\comm{Dans l’octave de l’Ascension}\liturgie{Intr. \textit{Ecce}}\nov{Neuvaine au Saint-Esprit, avant la Pentecôte — 6\textsuperscript{e} jour}{\pg{56}}}
\semainedroite{du 9 au 15 mai 2027}
\jourcal{Jeudi}{13}{0}{blanc}{\offgrand{Octave de l’Ascension}}{dbl maj.}{\comm{S. Robert Bellarmin, Évêque, Confesseur et Docteur de l’Église}\liturgie{Intr. \textit{Viri Galilǽi}}\nov{Neuvaine au Saint-Esprit, avant la Pentecôte — 7\textsuperscript{e} jour}{\pg{56}}}
\jourcal{Vendredi}{14}{0}{blanc}{\offpetit{Vendredi après l’octave de l’Ascension}}{semi-dbl}{\comm{S. Boniface, Martyr}\liturgie{Intr. \textit{Exáudi}}\jeune{Abstinence}\nov{Neuvaine au Saint-Esprit, avant la Pentecôte — 8\textsuperscript{e} jour}{\pg{56}}}
\jourcal{Samedi}{15}{0}{rouge}{\offpetit{Vigile de la Pentecôte}}{vigile 1\textsuperscript{re} cl.}{\comm{S. Jean-Baptiste de La Salle, Confesseur}\liturgie{Intr. \textit{Cum sanctificátus fúero in vobis} \enspace·\enspace \textit{Statio ad S. Joannem in Laterano}}\jeune{Jeûne et abstinence}\dev{Quinze samedis du Rosaire — 5/15}{\pg{74}}\nov{Neuvaine au Saint-Esprit, avant la Pentecôte — 9\textsuperscript{e} jour}{\pg{56}}}
\notes
\semainegauche{du 16 au 22 mai 2027}{Dimanche de la Pentecôte}
\jourcal{Dimanche}{16}{1}{rouge}{\offgrand{Dimanche de la Pentecôte}}{1\textsuperscript{re} cl.}{\comm{S. Ubald, Évêque et Confesseur}\liturgie{Intr. \textit{Spíritus Dómini replévit orbem terrárum} \enspace·\enspace \textit{Statio ad S. Petrum}}\dev{Six dimanches de S. Louis de Gonzague — 6/6}{\pg{63}}\dev{Veni Creator}{\pg{56}}}
\jourfete{Lundi}{17}{1}{rouge}{\offgrand{Lundi de Pentecôte}}{1\textsuperscript{re} cl.}{\liturgie{Intr. \textit{Cibávit eos ex ádipe fruménti} \enspace·\enspace \textit{Statio ad S. Petrum ad Vincula}}}
\jourfete{Mardi}{18}{1}{rouge}{\offgrand{Mardi de Pentecôte}}{1\textsuperscript{re} cl.}{\liturgie{Intr. \textit{Accípite jucunditátem glóriæ vestræ} \enspace·\enspace \textit{Statio ad S. Anastasiam}}\dev{Treize mardis de S. Antoine — 10/13}{\pg{68}}}
\jourfete{Mercredi}{19}{1}{rouge}{\offgrand{Mercredi des Quatre-Temps de Pentecôte}}{1\textsuperscript{re} cl.}{\liturgie{Intr. \textit{Deus} \enspace·\enspace \textit{Statio ad S. Mariam majorem}}\jeune{Jeûne et abstinence}}
\semainedroite{du 16 au 22 mai 2027}
\jourfete{Jeudi}{20}{1}{rouge}{\offgrand{Jeudi de Pentecôte}}{1\textsuperscript{re} cl.}{\liturgie{Intr. \textit{Spíritus Dómini replévit orbem terrárum} \enspace·\enspace \textit{Statio ad S. Laurentium extra muros}}}
\jourfete{Vendredi}{21}{1}{rouge}{\offgrand{Vendredi des Quatre-Temps de Pentecôte}}{1\textsuperscript{re} cl.}{\liturgie{Intr. \textit{Repleátur os meum laude tua} \enspace·\enspace \textit{Statio ad Ss. duodecim Apostolos}}\jeune{Jeûne et abstinence}}
\jourfete{Samedi}{22}{1}{rouge}{\offgrand{Samedi des Quatre-Temps de Pentecôte}}{1\textsuperscript{re} cl.}{\liturgie{Intr. \textit{Cáritas Dei diffúsa est in…} \enspace·\enspace \textit{Statio ad S. Petrum}}\jeune{Jeûne et abstinence}\dev{Quinze samedis du Rosaire — 6/15}{\pg{74}}}
\notes
\semainegauche{du 23 au 29 mai 2027}{Très Sainte Trinité}
\jourcal{Dimanche}{23}{1}{blanc}{\offgrand{Très Sainte Trinité}}{1\textsuperscript{re} cl.}{\liturgie{Intr. \textit{Benedícta sit sancta Trínitas atque…}}\dev{Dix dimanches de S. Ignace — 1/10}{\pg{66}}}
\jourcal{Lundi}{24}{0}{vert}{\offpetit{Férie}}{férie}{\liturgie{Intr. \textit{Benedícta sit sancta Trínitas atque…}}}
\jourcal{Mardi}{25}{0}{blanc}{\offpetit{S. Grégoire VII, Pape et Confesseur}}{dbl}{\comm{S. Urbain, Pape et Martyr}\liturgie{Intr. \textit{Státuit ei Dóminus testaméntum pacis}}\dev{Treize mardis de S. Antoine — 11/13}{\pg{68}}}
\jourcal{Mercredi}{26}{0}{blanc}{\offpetit{S. Philippe Neri, Confesseur}}{dbl}{\comm{S. Éleuthère, Pape et Martyr}\liturgie{Intr. \textit{Cáritas Dei diffúsa est in…}}\nov{Neuvaine au Sacré-Cœur — 1\textsuperscript{er} jour}{}}
\semainedroite{du 23 au 29 mai 2027}
\jourfete{Jeudi}{27}{1}{blanc}{\offgrand{Fête-Dieu (Saint-Sacrement)}}{1\textsuperscript{re} cl.}{\comm{S. Bède le Vénérable, Confesseur et Docteur de l’Église ; S. Jean I\textsuperscript{er}, Pape et Martyr}\liturgie{Intr. \textit{Cibávit eos ex ádipe fruménti}}\nov{Neuvaine au Sacré-Cœur — 2\textsuperscript{e} jour}{}}
\jourfete{Vendredi}{28}{1}{blanc}{\offgrand{Dans l’octave du Saint-Sacrement}}{2\textsuperscript{e} cl.}{\comm{S. Augustin de Cantorbéry, Évêque et Confesseur}\liturgie{Intr. \textit{Benedícta sit sancta Trínitas atque…}}\jeune{Abstinence}\nov{Neuvaine au Sacré-Cœur — 3\textsuperscript{e} jour}{}}
\jourfete{Samedi}{29}{1}{blanc}{\offgrand{Dans l’octave du Saint-Sacrement}}{2\textsuperscript{e} cl.}{\comm{Ste Marie-Madeleine de Pazzi, Vierge}\liturgie{Intr. \textit{Benedícta sit sancta Trínitas atque…}}\dev{Quinze samedis du Rosaire — 7/15}{\pg{74}}\nov{Neuvaine au Sacré-Cœur — 4\textsuperscript{e} jour}{}}
\notes
\moisgauche{Juin}{2027}{\devmois{Mois du Sacré-Cœur de Jésus}{253}{1}}{Les pensées de son Cœur sont de génération en génération, pour arracher à la mort leurs âmes et les nourrir dans leur faim.}{Ps 32:11; 32:19}{Sacré-Cœur}
\moisdroite{Juin 2027}{\rubmois{Fêtes principales}\lignemois{1}{Dans l’octave du Saint-Sacrement}\lignemois{2}{Dans l’octave du Saint-Sacrement}\lignemois{4}{Sacré-Cœur de Jésus}\lignemois{5}{Bienheureuse Vierge Marie, Reine}\lignemois{24}{Nativité de Saint Jean-Baptiste}\lignemois{29}{Les Ss. Apôtres Pierre et Paul}\lignemois{30}{En la Commémoraison de l’Apôtre S. Paul}\rubmois{Dévotions et neuvaines qui commencent}\lignemois{15}{Quinze mardis de S. Dominique\pg{68}}\rubmois{Pénitence}\lignemoislong{Abstinence}{4, 11, 18, 25}}
\semainegauche{du 30 mai au 5 juin 2027}{2\textsuperscript{e} dimanche après la Pentecôte}
\jourcal{Dimanche}{30}{1}{blanc}{\offgrand{2\textsuperscript{e} dimanche après la Pentecôte}}{dim. 2\textsuperscript{e} cl.}{\comm{S. Félix I\textsuperscript{er}, Pape et Martyr ; Octave du Saint-Sacrement}\liturgie{Intr. \textit{Factus est Dóminus protéctor meus}}\dev{Dix dimanches de S. Ignace — 2/10}{\pg{66}}\nov{Neuvaine au Sacré-Cœur — 5\textsuperscript{e} jour}{}}
\jourfete{Lundi}{31}{1}{blanc}{\offgrand{Dans l’octave du Saint-Sacrement}}{2\textsuperscript{e} cl.}{\comm{Ste Pétronille, Vierge}\liturgie{Intr. \textit{Cibávit eos ex ádipe fruménti}}\nov{Neuvaine au Sacré-Cœur — 6\textsuperscript{e} jour}{}}
\jourfete{Mardi}{1}{1}{blanc}{\offgrand{Dans l’octave du Saint-Sacrement}}{2\textsuperscript{e} cl.}{\comm{Ste Angèle Mérici, Vierge}\liturgie{Intr. \textit{Cibávit eos ex ádipe fruménti}}\dev{Treize mardis de S. Antoine — 12/13}{\pg{68}}\nov{Neuvaine au Sacré-Cœur — 7\textsuperscript{e} jour}{}}
\jourfete{Mercredi}{2}{1}{blanc}{\offgrand{Dans l’octave du Saint-Sacrement}}{2\textsuperscript{e} cl.}{\comm{Ss. Marcellin, Pierre et Érasme, Martyrs}\liturgie{Intr. \textit{Cibávit eos ex ádipe fruménti}}\dev{Premier mercredi : S. Joseph}{\pg{68}}\nov{Neuvaine au Sacré-Cœur — 8\textsuperscript{e} jour}{}}
\semainedroite{du 30 mai au 5 juin 2027}
\jourcal{Jeudi}{3}{0}{blanc}{\offgrand{Octave du Saint-Sacrement}}{dbl maj.}{\liturgie{Intr. \textit{Cibávit eos ex ádipe fruménti}}\dev{Heure sainte (veille du premier vendredi)}{\pg{69}}\nov{Neuvaine au Sacré-Cœur — 9\textsuperscript{e} jour}{}}
\jourfete{Vendredi}{4}{1}{blanc}{\offgrand{Sacré-Cœur de Jésus}}{1\textsuperscript{re} cl.}{\comm{S. François Caracciolo, Confesseur}\liturgie{Intr. \textit{Cogitatiónes Cordis ejus in generatióne…}}\jeune{Abstinence}\dev{Premier vendredi du mois — 7/9}{\pg{69}}}
\jourfete{Samedi}{5}{1}{blanc}{\offgrand{Bienheureuse Vierge Marie, Reine}}{2\textsuperscript{e} cl.}{\comm{Dans l’octave du Sacré-Cœur ; S. Boniface, Évêque et Martyr}\liturgie{Intr. \textit{Gaudeámus omnes in Dómino}}\noticesaint{Fête de la bienheureuse Vierge Marie, Reine.}\dev{Quinze samedis du Rosaire — 8/15}{\pg{74}}\dev{Premier samedi du mois}{\pg{73}}}
\notes
\semainegauche{du 6 au 12 juin 2027}{3\textsuperscript{e} dimanche après la Pentecôte}
\jourcal{Dimanche}{6}{1}{blanc}{\offgrand{3\textsuperscript{e} dimanche après la Pentecôte}}{dim.}{\comm{S. Norbert, Évêque et Confesseur ; Octave du Sacré-Cœur}\liturgie{Intr. \textit{Réspice in me et miserére…}}\dev{Dix dimanches de S. Ignace — 3/10}{\pg{66}}}
\jourcal{Lundi}{7}{0}{blanc}{\offpetit{Dans l’octave du Sacré-Cœur}}{semi-dbl}{\liturgie{Intr. \textit{Cogitatiónes Cordis ejus in generatióne…}}}
\jourcal{Mardi}{8}{0}{blanc}{\offpetit{Dans l’octave du Sacré-Cœur}}{semi-dbl}{\liturgie{Intr. \textit{Cogitatiónes Cordis ejus in generatióne…}}\dev{Treize mardis de S. Antoine — 13/13}{\pg{68}}}
\jourcal{Mercredi}{9}{0}{blanc}{\offpetit{Dans l’octave du Sacré-Cœur}}{semi-dbl}{\comm{Ss. Prime et Félicien, Martyrs}\liturgie{Intr. \textit{Cogitatiónes Cordis ejus in generatióne…}}}
\semainedroite{du 6 au 12 juin 2027}
\jourcal{Jeudi}{10}{0}{blanc}{\offpetit{Ste Marguerite, Reine et Veuve}}{semi-dbl}{\comm{Dans l’octave du Sacré-Cœur}\liturgie{Intr. \textit{Cognóvi}}}
\jourcal{Vendredi}{11}{0}{blanc}{\offgrand{Octave du Sacré-Cœur}}{dbl maj.}{\comm{S. Barnabé, Apôtre}\liturgie{Intr. \textit{Cogitatiónes Cordis ejus in generatióne…}}\jeune{Abstinence}}
\jourcal{Samedi}{12}{0}{blanc}{\offpetit{S. Jean de S. Facond, Confesseur}}{dbl}{\comm{Ss. Basilide, Cyrin, Nabor et Nazaire, Martyrs}\liturgie{Intr. \textit{Os justi meditábitur sapiéntiam}}\dev{Quinze samedis du Rosaire — 9/15}{\pg{74}}}
\notes
\semainegauche{du 13 au 19 juin 2027}{4\textsuperscript{e} dimanche après la Pentecôte}
\jourcal{Dimanche}{13}{1}{vert}{\offgrand{4\textsuperscript{e} dimanche après la Pentecôte}}{dim.}{\comm{S. Antoine de Padoue, Confesseur}\liturgie{Intr. \textit{Dóminus illuminátio mea et salus…}}\dev{Dix dimanches de S. Ignace — 4/10}{\pg{66}}}
\jourcal{Lundi}{14}{0}{blanc}{\offpetit{S. Basile le Grand, Confesseur et Docteur de l’Église}}{dbl}{\liturgie{Intr. \textit{In médio Ecclésiæ apéruit os…}}}
\jourcal{Mardi}{15}{0}{rouge}{\offpetit{Ss. Vite, Modeste et Crescence, Martyrs}}{simple}{\liturgie{Intr. \textit{Multæ tribulatiónes justórum}}\dev{Quinze mardis de S. Dominique — 1/15}{\pg{68}}}
\jourcal{Mercredi}{16}{0}{vert}{\offpetit{Férie}}{férie}{\liturgie{Intr. \textit{Dóminus illuminátio mea et salus…}}}
\semainedroite{du 13 au 19 juin 2027}
\jourcal{Jeudi}{17}{0}{vert}{\offpetit{Férie}}{férie}{\liturgie{Intr. \textit{Dóminus illuminátio mea et salus…}}}
\jourcal{Vendredi}{18}{0}{blanc}{\offpetit{S. Éphrem le Syrien, Confesseur et Docteur de l’Église}}{dbl}{\comm{Ss. Marc et Marcellien, Martyrs}\liturgie{Intr. \textit{In médio Ecclésiæ apéruit os…}}\jeune{Abstinence}}
\jourcal{Samedi}{19}{0}{blanc}{\offpetit{Ste Julienne de Falconieri, Vierge}}{dbl}{\comm{Ss. Gervais et Protais, Martyrs}\liturgie{Intr. \textit{Dilexísti justítiam}}\dev{Quinze samedis du Rosaire — 10/15}{\pg{74}}}
\notes
\semainegauche{du 20 au 26 juin 2027}{5\textsuperscript{e} dimanche après la Pentecôte}
\jourcal{Dimanche}{20}{1}{vert}{\offgrand{5\textsuperscript{e} dimanche après la Pentecôte}}{dim.}{\comm{S. Silvère, Pape et Martyr}\liturgie{Intr. \textit{Exáudi}}\dev{Dix dimanches de S. Ignace — 5/10}{\pg{66}}}
\jourcal{Lundi}{21}{0}{blanc}{\offpetit{S. Louis de Gonzague, Confesseur}}{dbl}{\liturgie{Intr. \textit{Minuísti eum paulo minus ab…}}}
\jourcal{Mardi}{22}{0}{blanc}{\offpetit{S. Paulin de Nole, Évêque et Confesseur}}{dbl}{\liturgie{Intr. \textit{Sacerdótes tui}}\dev{Quinze mardis de S. Dominique — 2/15}{\pg{68}}}
\jourcal{Mercredi}{23}{0}{violet}{\offpetit{Vigile de la Nativité de S. Jean-Baptiste}}{vigile}{\liturgie{Intr. \textit{Ne tímeas}}}
\semainedroite{du 20 au 26 juin 2027}
\jourfete{Jeudi}{24}{1}{blanc}{\offgrand{Nativité de Saint Jean-Baptiste}}{1\textsuperscript{re} cl.}{\liturgie{Intr. \textit{De ventre matris meæ vocávit…}}\noticesaint{La Nativité de saint Jean-Baptiste, précurseur du Seigneur, fils des saints Zacharie et Élisabeth. Il fut rempli du Saint-Esprit dès le sein de sa mère.}}
\jourcal{Vendredi}{25}{0}{blanc}{\offpetit{S. Guillaume, Abbé}}{dbl}{\comm{2\textsuperscript{e} jour dans l’Octave de la Nativité de S. Jean-Baptiste}\liturgie{Intr. \textit{Os justi meditábitur sapiéntiam}}\jeune{Abstinence}}
\jourcal{Samedi}{26}{0}{rouge}{\offpetit{Ss. Jean et Paul, Martyrs}}{dbl}{\comm{3\textsuperscript{e} jour dans l’octave de S. Jean-Baptiste}\liturgie{Intr. \textit{Multæ tribulatiónes justórum}}\dev{Quinze samedis du Rosaire — 11/15}{\pg{74}}}
\notes
\moisgauche{Juillet}{2027}{\devmois{Mois du Précieux Sang}{217}{1}}{Vous nous avez rachetés, Seigneur, par votre Sang, de toute tribu, de toute langue, de tout peuple et de toute nation : et vous nous avez fait royaume pour notre Dieu.}{Apoc 5:9-10}{Précieux Sang}
\moisdroite{Juillet 2027}{\rubmois{Fêtes principales}\lignemois{1}{Le Très Précieux Sang de Notre Seigneur Jésus-Christ}\lignemois{2}{Visitation de la Bienheureuse Vierge Marie}\lignemois{6}{Octave des Ss. Apôtres Pierre et Paul}\lignemois{16}{Notre-Dame du Mont-Carmel}\lignemois{25}{S. Jacques, Apôtre}\lignemois{26}{Ste Anne, Mère de la Bienheureuse Vierge Marie}\lignemois{31}{S. Ignace, Confesseur}\rubmois{Dévotions et neuvaines qui commencent}\lignemois{30}{Sept vendredis des Douleurs\pg{72}}\lignemois{31}{Sept samedis de N.-D. de la Merci\pg{76}}\rubmois{Pénitence}\lignemoislong{Abstinence}{2, 9, 16, 23, 30}}
\semainegauche{du 27 juin au 3 juillet 2027}{6\textsuperscript{e} dimanche après la Pentecôte}
\jourcal{Dimanche}{27}{1}{vert}{\offgrand{6\textsuperscript{e} dimanche après la Pentecôte}}{dim.}{\comm{4\textsuperscript{e} jour dans l’Octave de la Nativité de S. Jean-Baptiste}\liturgie{Intr. \textit{Dóminus fortitudo plebis suæ}}\dev{Dix dimanches de S. Ignace — 6/10}{\pg{66}}}
\jourcal{Lundi}{28}{0}{rouge}{\offpetit{S. Irénée, Évêque et Martyr}}{dbl}{\comm{Vigile des Ss. Pierre et Paul, Apôtres ; 5\textsuperscript{e} jour dans l’octave de S. Jean-Baptiste}\liturgie{Intr. \textit{Lex veritátis fuit in ore…}}}
\jourfete{Mardi}{29}{1}{rouge}{\offgrand{Les Ss. Apôtres Pierre et Paul}}{1\textsuperscript{re} cl.}{\liturgie{Intr. \textit{Nunc scio vere}}\noticesaint{À Rome, l’anniversaire des saints Apôtres Pierre et Paul, qui souffrirent la même année et le même jour sous l’empereur Néron. […]}\dev{Quinze mardis de S. Dominique — 3/15}{\pg{68}}}
\jourcal{Mercredi}{30}{0}{rouge}{\offgrand{En la Commémoraison de l’Apôtre S. Paul}}{dbl maj.}{\comm{7\textsuperscript{e} jour dans l’octave de S. Jean-Baptiste}\liturgie{Intr. \textit{Scio}}\noticesaint{La Commémoraison de l’Apôtre saint Paul.}}
\semainedroite{du 27 juin au 3 juillet 2027}
\jourfete{Jeudi}{1}{1}{rouge}{\offgrand{Le Très Précieux Sang de Notre Seigneur Jésus-Christ}}{1\textsuperscript{re} cl.}{\comm{Octave de S. Jean-Baptiste}\liturgie{Intr. \textit{Redemísti nos}}\noticesaint{La fête du Très Précieux Sang de Notre Seigneur Jésus-Christ.}\dev{Heure sainte (veille du premier vendredi)}{\pg{69}}}
\jourfete{Vendredi}{2}{1}{blanc}{\offgrand{Visitation de la Bienheureuse Vierge Marie}}{2\textsuperscript{e} cl.}{\comm{Ss. Processus et Martinien, Martyrs}\noticesaint{La Visitation de la bienheureuse Vierge Marie chez Élisabeth.}\jeune{Abstinence}\dev{Premier vendredi du mois — 8/9}{\pg{69}}}
\jourcal{Samedi}{3}{0}{blanc}{\offpetit{S. Léon, Pape et Confesseur}}{semi-dbl}{\comm{5\textsuperscript{e} jour dans l’Octave des Ss. Apôtres Pierre et Paul}\liturgie{Intr. \textit{Státuit ei Dóminus testaméntum pacis}}\dev{Quinze samedis du Rosaire — 12/15}{\pg{74}}\dev{Premier samedi du mois}{\pg{73}}}
\notes
\semainegauche{du 4 au 10 juillet 2027}{7\textsuperscript{e} dimanche après la Pentecôte}
\jourcal{Dimanche}{4}{1}{vert}{\offgrand{7\textsuperscript{e} dimanche après la Pentecôte}}{dim.}{\comm{6\textsuperscript{e} jour dans l’Octave des Ss. Apôtres Pierre et Paul}\liturgie{Intr. \textit{Omnes gentes}}\dev{Dix dimanches de S. Ignace — 7/10}{\pg{66}}}
\jourcal{Lundi}{5}{0}{blanc}{\offpetit{S. Antoine Marie Zaccaria, Confesseur}}{dbl}{\comm{7\textsuperscript{e} jour dans l’Octave des Ss. Apôtres Pierre et Paul}\liturgie{Intr. \textit{Sermo meus et prædicátio mea…}}}
\jourcal{Mardi}{6}{0}{rouge}{\offgrand{Octave des Ss. Apôtres Pierre et Paul}}{dbl maj.}{\liturgie{Intr. \textit{Sapiéntiam Sanctórum narrent populi}}\dev{Quinze mardis de S. Dominique — 4/15}{\pg{68}}}
\jourcal{Mercredi}{7}{0}{blanc}{\offpetit{Ss. Cyrille et Méthode, Évêques et Confesseurs}}{dbl}{\liturgie{Intr. \textit{Sacerdótes tui}}\dev{Premier mercredi : S. Joseph}{\pg{68}}}
\semainedroite{du 4 au 10 juillet 2027}
\jourcal{Jeudi}{8}{0}{blanc}{\offpetit{Ste Élisabeth, Reine du Portugal, Veuve}}{semi-dbl}{\liturgie{Intr. \textit{Cognóvi}}}
\jourcal{Vendredi}{9}{0}{vert}{\offpetit{Férie}}{férie}{\liturgie{Intr. \textit{Omnes gentes}}\jeune{Abstinence}}
\jourcal{Samedi}{10}{0}{rouge}{\offpetit{Les Sept Frères Martyrs, et Stes Rufine et Seconde, Vierges et Martyres}}{semi-dbl}{\liturgie{Intr. \textit{Laudáte}}\dev{Quinze samedis du Rosaire — 13/15}{\pg{74}}}
\notes
\semainegauche{du 11 au 17 juillet 2027}{8\textsuperscript{e} dimanche après la Pentecôte}
\jourcal{Dimanche}{11}{1}{vert}{\offgrand{8\textsuperscript{e} dimanche après la Pentecôte}}{dim.}{\comm{S. Pie I\textsuperscript{er}, Pape et Martyr}\liturgie{Intr. \textit{Suscépimus}}\dev{Dix dimanches de S. Ignace — 8/10}{\pg{66}}}
\jourcal{Lundi}{12}{0}{blanc}{\offpetit{S. Jean Gualbert, Abbé}}{dbl}{\comm{Ss. Nabor et Félix, Martyrs}\liturgie{Intr. \textit{Os justi meditábitur sapiéntiam}}}
\jourcal{Mardi}{13}{0}{rouge}{\offpetit{S. Anaclet, Pape et Martyr}}{semi-dbl}{\liturgie{Intr. \textit{Sacerdótes tui}}\dev{Quinze mardis de S. Dominique — 5/15}{\pg{68}}}
\jourcal{Mercredi}{14}{0}{blanc}{\offpetit{S. Bonaventure, Évêque, Confesseur et Docteur de l’Église}}{dbl}{\liturgie{Intr. \textit{In médio Ecclésiæ apéruit os…}}}
\semainedroite{du 11 au 17 juillet 2027}
\jourcal{Jeudi}{15}{0}{blanc}{\offpetit{S. Henri, Empereur et Confesseur}}{semi-dbl}{\liturgie{Intr. \textit{Os justi meditábitur sapiéntiam}}}
\jourcal{Vendredi}{16}{0}{blanc}{\offgrand{Notre-Dame du Mont-Carmel}}{dbl maj.}{\liturgie{Intr. \textit{Gaudeamus omnes in Domino diem…}}\jeune{Abstinence}}
\jourcal{Samedi}{17}{0}{blanc}{\offpetit{S. Alexis, Confesseur}}{semi-dbl}{\liturgie{Intr. \textit{Os justi meditábitur sapiéntiam}}\dev{Quinze samedis du Rosaire — 14/15}{\pg{74}}}
\notes
\semainegauche{du 18 au 24 juillet 2027}{9\textsuperscript{e} dimanche après la Pentecôte}
\jourcal{Dimanche}{18}{1}{vert}{\offgrand{9\textsuperscript{e} dimanche après la Pentecôte}}{dim.}{\comm{S. Camille de Lellis, Confesseur ; Ste Symphorose et ses sept Fils, Martyrs}\liturgie{Intr. \textit{Ecce}}\dev{Dix dimanches de S. Ignace — 9/10}{\pg{66}}}
\jourcal{Lundi}{19}{0}{blanc}{\offpetit{S. Vincent de Paul, Confesseur}}{dbl}{\liturgie{Intr. \textit{Justus ut palma florébit}}}
\jourcal{Mardi}{20}{0}{blanc}{\offpetit{S. Jérôme Émilien, Confesseur}}{dbl}{\comm{Ste Marguerite, Vierge et Martyre}\liturgie{Intr. \textit{Effúsum est in terra jecur…}}\dev{Quinze mardis de S. Dominique — 6/15}{\pg{68}}}
\jourcal{Mercredi}{21}{0}{blanc}{\offpetit{Ste Praxède, Vierge}}{simple}{\liturgie{Intr. \textit{Loquébar de testimóniis tuis in…}}}
\semainedroite{du 18 au 24 juillet 2027}
\jourcal{Jeudi}{22}{0}{blanc}{\offpetit{Ste Marie-Madeleine, pénitente}}{dbl}{\liturgie{Intr. \textit{Me exspectavérunt peccatóres}}}
\jourcal{Vendredi}{23}{0}{rouge}{\offpetit{S. Apollinaire, Évêque et Martyr}}{dbl}{\comm{S. Liboire, Évêque et Confesseur}\liturgie{Intr. \textit{Sacerdótes Dei}}\jeune{Abstinence}}
\jourcal{Samedi}{24}{0}{violet}{\offpetit{Vigile de S. Jacques, Apôtre}}{vigile}{\comm{Ste Christine, Vierge et Martyre}\liturgie{Intr. \textit{Ego autem}}\dev{Quinze samedis du Rosaire — 15/15}{\pg{74}}}
\notes
\semainegauche{du 25 au 31 juillet 2027}{10\textsuperscript{e} dimanche après la Pentecôte}
\jourfete{Dimanche}{25}{1}{rouge}{\offgrand{S. Jacques, Apôtre}}{2\textsuperscript{e} cl.}{\comm{10\textsuperscript{e} dimanche après la Pentecôte ; S. Christophe, Martyr}\liturgie{Intr. \textit{Mihi autem nimis honoráti sunt…}}\noticesaint{Saint Jacques Apôtre, frère de l’Évangéliste saint Jean. Décapité par ordre d’Hérode Agrippa, au moment de la fête de Pâques, le premier parmi les Apôtres il reçut la couronne du martyre. […]}\dev{Dix dimanches de S. Ignace — 10/10}{\pg{66}}}
\jourfete{Lundi}{26}{1}{blanc}{\offgrand{Ste Anne, Mère de la Bienheureuse Vierge Marie}}{2\textsuperscript{e} cl.}{\liturgie{Intr. \textit{Gaudeámus omnes in Dómino}}\noticesaint{Le paisible sommeil de sainte Anne, mère de l’Immaculée Vierge Marie, Mère de Dieu.}}
\jourcal{Mardi}{27}{0}{rouge}{\offpetit{S. Pantaléon, Martyr}}{simple}{\liturgie{Intr. \textit{Lætábitur justus in Dómino}}\dev{Quinze mardis de S. Dominique — 7/15}{\pg{68}}}
\jourcal{Mercredi}{28}{0}{rouge}{\offpetit{Ss. Nazaire et Celse, Martyrs, Victor, Pape et Martyr et Innocent I\textsuperscript{er}, Pape et Confesseur}}{semi-dbl}{\liturgie{Intr. \textit{Intret in conspéctu tuo}}}
\semainedroite{du 25 au 31 juillet 2027}
\jourcal{Jeudi}{29}{0}{blanc}{\offpetit{Ste Marthe, Vierge}}{semi-dbl}{\comm{S. Félix, Pape, et les Ss. Simplice, Faustin et Béatrice, Martyrs}\liturgie{Intr. \textit{Dilexísti justítiam}}}
\jourcal{Vendredi}{30}{0}{rouge}{\offpetit{Ss. Abdon et Sennen, Martyrs}}{simple}{\liturgie{Intr. \textit{Intret in conspéctu tuo}}\jeune{Abstinence}\dev{Sept vendredis des Douleurs — 1/7}{\pg{72}}}
\jourcal{Samedi}{31}{0}{blanc}{\offgrand{S. Ignace, Confesseur}}{dbl maj.}{\liturgie{Intr. \textit{In nómine Jesu omne genu…}}\noticesaint{À Rome, l’anniversaire de saint Ignace, Prêtre et Confesseur, fondateur de la Compagnie de Jésus. […]}\dev{Sept samedis de N.-D. de la Merci — 1/7}{\pg{76}}}
\notes
\moisgauche{Août}{2027}{\devmois{Mois du Cœur Immaculé de Marie}{389}{1}}{Il parut dans le ciel un grand signe : une femme revêtue du soleil, la lune sous ses pieds, et une couronne de douze étoiles sur sa tête.}{Ap 12:1}{Assomption}
\moisdroite{Août 2027}{\rubmois{Fêtes principales}\lignemois{4}{S. Dominique, Confesseur}\lignemois{5}{Dédicace de Sainte-Marie-aux-Neiges}\lignemois{6}{Transfiguration de Notre-Seigneur}\lignemois{10}{S. Laurent, Martyr}\lignemois{15}{Assomption de la Bienheureuse Vierge Marie}\lignemois{16}{S. Joachim, Père de la B. V. Marie}\lignemois{22}{Cœur Immaculé de Marie}\lignemois{24}{S. Barthélemy, Apôtre}\rubmois{Dévotions et neuvaines qui commencent}\lignemois{6}{Neuvaine de l’Assomption}\lignemois{15}{Cinq dimanches des Stigmates de S. François\pg{62}}\lignemois{30}{Neuvaine de la Nativité de Notre-Dame}\rubmois{Pénitence}\lignemoislong{Jeûne et abstinence}{14}\lignemoislong{Abstinence}{6, 13, 20, 27}}
\semainegauche{du 1 au 7 août 2027}{11\textsuperscript{e} dimanche après la Pentecôte}
\jourcal{Dimanche}{1}{1}{vert}{\offgrand{11\textsuperscript{e} dimanche après la Pentecôte}}{dim.}{\comm{S. Pierre aux Liens ; Ss. Macchabées, Martyrs}\liturgie{Intr. \textit{Deus in loco sancto suo}}}
\jourcal{Lundi}{2}{0}{blanc}{\offpetit{S. Alphonse Marie de Liguori, Évêque, Confesseur et Docteur de l’Église}}{dbl}{\comm{S. Étienne I\textsuperscript{er}, Pape et Martyr}\liturgie{Intr. \textit{Spíritus Dómini super me}}}
\jourcal{Mardi}{3}{0}{rouge}{\offpetit{Invention de S. Étienne, premier, Martyr}}{semi-dbl}{\liturgie{Intr. \textit{Sedérunt príncipes}}\dev{Quinze mardis de S. Dominique — 8/15}{\pg{68}}}
\jourcal{Mercredi}{4}{0}{blanc}{\offgrand{S. Dominique, Confesseur}}{dbl maj.}{\liturgie{Intr. \textit{Os justi meditábitur sapiéntiam}}\noticesaint{Saint Dominique, Confesseur, fondateur de l’Ordre des Frères Prêcheurs. Il s’endormit dans la paix le 8 des ides de ce mois (6 août).}\dev{Premier mercredi : S. Joseph}{\pg{68}}}
\semainedroite{du 1 au 7 août 2027}
\jourcal{Jeudi}{5}{0}{blanc}{\offgrand{Dédicace de Sainte-Marie-aux-Neiges}}{dbl maj.}{\liturgie{Intr. \textit{Salve}}\noticesaint{À Rome, sur le Mont Esquilin, la dédicace de la basilique de Sainte-Marie-aux-Neiges.}\dev{Heure sainte (veille du premier vendredi)}{\pg{69}}}
\jourfete{Vendredi}{6}{1}{blanc}{\offgrand{Transfiguration de Notre-Seigneur}}{2\textsuperscript{e} cl.}{\comm{Ss. Sixte II, Pape, Félicissime et Agapit, Martyrs}\liturgie{Intr. \textit{Illuxérunt coruscatiónes tuæ orbi terræ}}\noticesaint{Sur le Mont Thabor, la Transfiguration de Notre Seigneur Jésus-Christ.}\jeune{Abstinence}\dev{Sept vendredis des Douleurs — 2/7}{\pg{72}}\dev{Premier vendredi du mois — 9/9}{\pg{69}}\nov{Neuvaine de l’Assomption — 1\textsuperscript{er} jour}{}}
\jourcal{Samedi}{7}{0}{blanc}{\offpetit{S. Gaétan de Thiène, Confesseur}}{dbl}{\comm{S. Donat, Évêque et Martyr}\liturgie{Intr. \textit{Os justi meditábitur sapiéntiam}}\dev{Sept samedis de N.-D. de la Merci — 2/7}{\pg{76}}\dev{Premier samedi du mois}{\pg{73}}\nov{Neuvaine de l’Assomption — 2\textsuperscript{e} jour}{}}
\notes
\semainegauche{du 8 au 14 août 2027}{12\textsuperscript{e} dimanche après la Pentecôte}
\jourcal{Dimanche}{8}{1}{vert}{\offgrand{12\textsuperscript{e} dimanche après la Pentecôte}}{dim.}{\comm{Ss. Cyriaque, Large et Smaragde, Martyrs}\liturgie{Intr. \textit{Deus}}\nov{Neuvaine de l’Assomption — 3\textsuperscript{e} jour}{}}
\jourcal{Lundi}{9}{0}{blanc}{\offpetit{S. Jean-Marie Vianney, Confesseur}}{dbl}{\comm{Vigile de S. Laurent, Martyr ; S. Romain, Martyr}\liturgie{Intr. \textit{Os justi meditábitur sapiéntiam}}\nov{Neuvaine de l’Assomption — 4\textsuperscript{e} jour}{}}
\jourfete{Mardi}{10}{1}{rouge}{\offgrand{S. Laurent, Martyr}}{2\textsuperscript{e} cl.}{\liturgie{Intr. \textit{Conféssio et pulchritúdo in conspéctu…}}\noticesaint{À Rome, sur la voie Tiburtine, l’anniversaire du bienheureux Laurent, Archidiacre. […]}\dev{Quinze mardis de S. Dominique — 9/15}{\pg{68}}\nov{Neuvaine de l’Assomption — 5\textsuperscript{e} jour}{}}
\jourcal{Mercredi}{11}{0}{rouge}{\offpetit{S. Tiburce et Ste Suzanne, Vierge, Martyrs}}{simple}{\liturgie{Intr. \textit{Salus autem justórum a Dómino}}\nov{Neuvaine de l’Assomption — 6\textsuperscript{e} jour}{}}
\semainedroite{du 8 au 14 août 2027}
\jourcal{Jeudi}{12}{0}{blanc}{\offpetit{Ste Claire, Vierge}}{dbl}{\liturgie{Intr. \textit{Dilexísti justítiam}}\nov{Neuvaine de l’Assomption — 7\textsuperscript{e} jour}{}}
\jourcal{Vendredi}{13}{0}{rouge}{\offpetit{Ss. Hippolyte et Cassien, Martyrs}}{simple}{\liturgie{Intr. \textit{Salus autem justórum a Dómino}}\jeune{Abstinence}\dev{Sept vendredis des Douleurs — 3/7}{\pg{72}}\nov{Neuvaine de l’Assomption — 8\textsuperscript{e} jour}{}}
\jourcal{Samedi}{14}{0}{violet}{\offpetit{Vigile de l’Assomption}}{vigile}{\comm{S. Eusèbe, Confesseur}\liturgie{Intr. \textit{Vultum tuum deprecabúntur omnes dívites…}}\jeune{Jeûne et abstinence}\dev{Sept samedis de N.-D. de la Merci — 3/7}{\pg{76}}\nov{Neuvaine de l’Assomption — 9\textsuperscript{e} jour}{}}
\notes
\semainegauche{du 15 au 21 août 2027}{13\textsuperscript{e} dimanche après la Pentecôte}
\jourfete{Dimanche}{15}{1}{blanc}{\offgrand{Assomption de la Bienheureuse Vierge Marie}}{1\textsuperscript{re} cl.}{\comm{13\textsuperscript{e} dimanche après la Pentecôte}\liturgie{Intr. \textit{Signum magnum appáruit in cœlo}}\noticesaint{L’Assomption de la très sainte Vierge Marie, Mère de Dieu.}\dev{Cinq dimanches des Stigmates de S. François — 1/5}{\pg{62}}}
\jourfete{Lundi}{16}{1}{blanc}{\offgrand{S. Joachim, Père de la B. V. Marie}}{2\textsuperscript{e} cl.}{\liturgie{Intr. \textit{Dispérsit}}\noticesaint{Saint Joachim, Confesseur, père de l’Immaculée Vierge Marie, Mère de Dieu. Son anniversaire est mentionné le 13 des calendes d’avril (20 mars).}}
\jourcal{Mardi}{17}{0}{blanc}{\offpetit{S. Hyacinthe, Confesseur}}{dbl}{\comm{3\textsuperscript{e} jour dans l’octave de l’Assomption ; Octave de S. Laurent, Martyr}\liturgie{Intr. \textit{Os justi meditábitur sapiéntiam}}\dev{Quinze mardis de S. Dominique — 10/15}{\pg{68}}}
\jourcal{Mercredi}{18}{0}{blanc}{\offpetit{4\textsuperscript{e} jour dans l’octave de l’Assomption}}{semi-dbl}{\comm{S. Agapit, Martyr}\liturgie{Intr. \textit{Signum magnum appáruit in cœlo}}}
\semainedroite{du 15 au 21 août 2027}
\jourcal{Jeudi}{19}{0}{blanc}{\offpetit{S. Jean Eudes, Confesseur}}{dbl}{\comm{5\textsuperscript{e} jour dans l’octave de l’Assomption}\liturgie{Intr. \textit{Os justi meditábitur sapiéntiam}}}
\jourcal{Vendredi}{20}{0}{blanc}{\offpetit{S. Bernard, Abbé et Docteur de l’Église}}{dbl}{\comm{6\textsuperscript{e} jour dans l’octave de l’Assomption}\liturgie{Intr. \textit{In médio Ecclésiæ apéruit os…}}\jeune{Abstinence}\dev{Sept vendredis des Douleurs — 4/7}{\pg{72}}}
\jourcal{Samedi}{21}{0}{blanc}{\offpetit{Ste Jeanne-Françoise Frémiot de Chantal, Veuve}}{dbl}{\comm{7\textsuperscript{e} jour dans l’octave de l’Assomption}\liturgie{Intr. \textit{Cognóvi}}\dev{Sept samedis de N.-D. de la Merci — 4/7}{\pg{76}}}
\notes
\semainegauche{du 22 au 28 août 2027}{14\textsuperscript{e} dimanche après la Pentecôte}
\jourfete{Dimanche}{22}{1}{blanc}{\offgrand{Cœur Immaculé de Marie}}{2\textsuperscript{e} cl.}{\comm{14\textsuperscript{e} dimanche après la Pentecôte ; Ss. Timothée, Hippolyte et Symphorien, Martyrs}\liturgie{Intr. \textit{Adeámus cum fidúcia ad thronum…}}\noticesaint{Fête du Cœur Immaculé de la bienheureuse Vierge Marie.}\dev{Cinq dimanches des Stigmates de S. François — 2/5}{\pg{62}}}
\jourcal{Lundi}{23}{0}{blanc}{\offpetit{S. Philippe Benizi, Confesseur}}{dbl}{\comm{Vigile de S. Barthélemy, Apôtre}\liturgie{Intr. \textit{Justus ut palma florébit}}}
\jourfete{Mardi}{24}{1}{rouge}{\offgrand{S. Barthélemy, Apôtre}}{2\textsuperscript{e} cl.}{\liturgie{Intr. \textit{Mihi autem nimis honoráti sunt…}}\noticesaint{Saint Barthélemy, Apôtre. Il prêcha l’Évangile du Christ dans les Indes. […]}\dev{Quinze mardis de S. Dominique — 11/15}{\pg{68}}}
\jourcal{Mercredi}{25}{0}{blanc}{\offpetit{S. Louis, Roi et Confesseur}}{semi-dbl}{\liturgie{Intr. \textit{Os justi meditábitur sapiéntiam}}}
\semainedroite{du 22 au 28 août 2027}
\jourcal{Jeudi}{26}{0}{rouge}{\offpetit{S. Zéphyrin, Pape et Martyr}}{simple}{}
\jourcal{Vendredi}{27}{0}{blanc}{\offpetit{S. Joseph Calasanz, Confesseur}}{dbl}{\liturgie{Intr. \textit{Veníte}}\jeune{Abstinence}\dev{Sept vendredis des Douleurs — 5/7}{\pg{72}}}
\jourcal{Samedi}{28}{0}{blanc}{\offpetit{S. Augustin, Évêque, Confesseur et Docteur de l’Église}}{dbl}{\comm{S. Hermès, Martyr}\liturgie{Intr. \textit{In médio Ecclésiæ apéruit os…}}\dev{Sept samedis de N.-D. de la Merci — 5/7}{\pg{76}}}
\notes
\moisgauche{Septembre}{2027}{\devmois{Mois de Notre-Dame des Sept Douleurs}{381}{1}}{Salut, ô Mère sainte ; mère qui avez enfanté le Roi qui régit le ciel et la terre dans les siècles des siècles.}{Sedulius}{Nativité de Notre-Dame}
\moisdroite{Septembre 2027}{\rubmois{Fêtes principales}\lignemois{8}{Nativité de la Bienheureuse Vierge Marie}\lignemois{14}{Exaltation de la Sainte Croix}\lignemois{15}{Les Sept Douleurs de la Bienheureuse Vierge Marie}\lignemois{21}{S. Matthieu, Apôtre et Évangéliste}\lignemois{24}{Notre-Dame de la Merci}\lignemois{29}{Dédicace de S. Michel Archange}\rubmois{Dévotions et neuvaines qui commencent}\lignemois{18}{Douze samedis avant l’Immaculée Conception\pg{73}}\lignemois{20}{Neuvaine à S. Michel\pg{28}}\lignemois{28}{Neuvaine du Rosaire}\rubmois{Pénitence}\lignemoislong{Jeûne et abstinence}{15, 17, 18}\lignemoislong{Abstinence}{3, 10, 24}}
\semainegauche{du 29 août au 4 septembre 2027}{15\textsuperscript{e} dimanche après la Pentecôte}
\jourcal{Dimanche}{29}{1}{vert}{\offgrand{15\textsuperscript{e} dimanche après la Pentecôte}}{dim.}{\comm{Décollation de S. Jean-Baptiste ; Ste Sabine, Martyre}\liturgie{Intr. \textit{Inclína}}\dev{Cinq dimanches des Stigmates de S. François — 3/5}{\pg{62}}}
\jourcal{Lundi}{30}{0}{blanc}{\offpetit{Ste Rose de Lima, Vierge}}{dbl}{\comm{Ss. Félix et Adaucte, Martyrs}\liturgie{Intr. \textit{Dilexísti justítiam}}\nov{Neuvaine de la Nativité de Notre-Dame — 1\textsuperscript{er} jour}{}}
\jourcal{Mardi}{31}{0}{blanc}{\offpetit{S. Raymond Nonnat, Confesseur}}{dbl}{\liturgie{Intr. \textit{Os justi meditábitur sapiéntiam}}\dev{Quinze mardis de S. Dominique — 12/15}{\pg{68}}\nov{Neuvaine de la Nativité de Notre-Dame — 2\textsuperscript{e} jour}{}}
\jourcal{Mercredi}{1}{0}{blanc}{\offpetit{S. Gilles, Abbé}}{simple}{\comm{Les Douze Saints Frères Martyrs}\liturgie{Intr. \textit{Os justi meditábitur sapiéntiam}}\dev{Premier mercredi : S. Joseph}{\pg{68}}\nov{Neuvaine de la Nativité de Notre-Dame — 3\textsuperscript{e} jour}{}}
\semainedroite{du 29 août au 4 septembre 2027}
\jourcal{Jeudi}{2}{0}{rouge}{\offpetit{S. Étienne, Roi et Confesseur}}{semi-dbl}{\liturgie{Intr. \textit{Os justi meditábitur sapiéntiam}}\dev{Heure sainte (veille du premier vendredi)}{\pg{69}}\nov{Neuvaine de la Nativité de Notre-Dame — 4\textsuperscript{e} jour}{}}
\jourcal{Vendredi}{3}{0}{blanc}{\offpetit{S. Pie X, Pape et Confesseur}}{dbl}{\liturgie{Intr. \textit{Éxtuli eléctum de pópulo}}\jeune{Abstinence}\dev{Sept vendredis des Douleurs — 6/7}{\pg{72}}\dev{Premier vendredi du mois}{\pg{69}}\nov{Neuvaine de la Nativité de Notre-Dame — 5\textsuperscript{e} jour}{}}
\jourcal{Samedi}{4}{0}{blanc}{\offpetit{Office de la Sainte Vierge le samedi}}{simple}{\liturgie{Intr. \textit{Salve, sancta Parens}}\dev{Sept samedis de N.-D. de la Merci — 6/7}{\pg{76}}\dev{Premier samedi du mois}{\pg{73}}\nov{Neuvaine de la Nativité de Notre-Dame — 6\textsuperscript{e} jour}{}}
\notes
\semainegauche{du 5 au 11 septembre 2027}{16\textsuperscript{e} dimanche après la Pentecôte}
\jourcal{Dimanche}{5}{1}{vert}{\offgrand{16\textsuperscript{e} dimanche après la Pentecôte}}{dim.}{\comm{S. Laurent Justinien, Évêque et Confesseur}\liturgie{Intr. \textit{Miserére mihi}}\dev{Cinq dimanches des Stigmates de S. François — 4/5}{\pg{62}}\nov{Neuvaine de la Nativité de Notre-Dame — 7\textsuperscript{e} jour}{}}
\jourcal{Lundi}{6}{0}{vert}{\offpetit{Férie}}{férie}{\liturgie{Intr. \textit{Miserére mihi}}\nov{Neuvaine de la Nativité de Notre-Dame — 8\textsuperscript{e} jour}{}}
\jourcal{Mardi}{7}{0}{vert}{\offpetit{Férie}}{férie}{\liturgie{Intr. \textit{Miserére mihi}}\dev{Quinze mardis de S. Dominique — 13/15}{\pg{68}}\nov{Neuvaine de la Nativité de Notre-Dame — 9\textsuperscript{e} jour}{}}
\jourfete{Mercredi}{8}{1}{blanc}{\offgrand{Nativité de la Bienheureuse Vierge Marie}}{2\textsuperscript{e} cl.}{\comm{S. Adrien, Martyr}\liturgie{Intr. \textit{Salve}}\noticesaint{La Nativité de la bienheureuse Marie, toujours Vierge, Mère de Dieu.}}
\semainedroite{du 5 au 11 septembre 2027}
\jourcal{Jeudi}{9}{0}{rouge}{\offpetit{S. Gorgon, Martyr}}{simple}{\liturgie{Intr. \textit{Lætábitur justus in Dómino}}}
\jourcal{Vendredi}{10}{0}{blanc}{\offpetit{S. Nicolas de Tolentino, Confesseur}}{dbl}{\liturgie{Intr. \textit{Justus ut palma florébit}}\jeune{Abstinence}\dev{Sept vendredis des Douleurs — 7/7}{\pg{72}}}
\jourcal{Samedi}{11}{0}{blanc}{\offpetit{Samedi dans l’octave de la Nativité de la Sainte Vierge}}{simple}{\comm{Ss. Prote et Hyacinthe, Martyrs}\liturgie{Intr. \textit{Salve}}\dev{Sept samedis de N.-D. de la Merci — 7/7}{\pg{76}}}
\notes
\semainegauche{du 12 au 18 septembre 2027}{17\textsuperscript{e} dimanche après la Pentecôte}
\jourcal{Dimanche}{12}{1}{vert}{\offgrand{17\textsuperscript{e} dimanche après la Pentecôte}}{dim.}{\comm{Saint Nom de Marie}\liturgie{Intr. \textit{Justus es}}\dev{Cinq dimanches des Stigmates de S. François — 5/5}{\pg{62}}}
\jourcal{Lundi}{13}{0}{vert}{\offpetit{Férie}}{férie}{\liturgie{Intr. \textit{Justus es}}}
\jourcal{Mardi}{14}{0}{rouge}{\offgrand{Exaltation de la Sainte Croix}}{dbl maj.}{\liturgie{Intr. \textit{Nos autem gloriári opórtet in…}}\noticesaint{L’Exaltation de la Sainte Croix, quand l’empereur Héraclius, vainqueur du roi Chosroès, la reporta de Perse à Jérusalem.}\dev{Quinze mardis de S. Dominique — 14/15}{\pg{68}}}
\jourfete{Mercredi}{15}{1}{blanc}{\offgrand{Les Sept Douleurs de la Bienheureuse Vierge Marie}}{2\textsuperscript{e} cl.}{\comm{Mercredi des Quatre-Temps de septembre ; S. Nicomède, Martyr}\liturgie{Intr. \textit{Stabant juxta Crucem Jesu Mater…} \enspace·\enspace \textit{Statio ad S. Mariam majorem}}\noticesaint{La fête des Sept Douleurs de la bienheureuse Vierge Marie.}\jeune{Jeûne et abstinence}\dev{Les sept Douleurs de Marie}{\pg{32}}}
\semainedroite{du 12 au 18 septembre 2027}
\jourcal{Jeudi}{16}{0}{rouge}{\offpetit{Ss. Corneille, Pape, et Cyprien, Évêque, Martyrs}}{semi-dbl}{\comm{Ss. Euphémie, Vierge, Lucie et Géminien, Martyrs}\liturgie{Intr. \textit{Intret in conspéctu tuo}}}
\jourcal{Vendredi}{17}{0}{blanc}{\offpetit{Impression des Stigmates de Saint François}}{dbl}{\comm{Vendredi des Quatre-Temps de septembre}\liturgie{Intr. \textit{Mihi autem absit gloriári} \enspace·\enspace \textit{Statio ad Ss. duodecim Apostolos}}\jeune{Jeûne et abstinence}}
\jourcal{Samedi}{18}{0}{blanc}{\offpetit{S. Joseph de Cupertino, Confesseur}}{dbl}{\comm{Samedi des Quatre-Temps de septembre}\liturgie{Intr. \textit{Diléctio Dei honorábilis sapiéntia} \enspace·\enspace \textit{Statio ad S. Petrum}}\jeune{Jeûne et abstinence}\dev{Douze samedis avant l’Immaculée Conception — 1/12}{\pg{73}}}
\notes
\semainegauche{du 19 au 25 septembre 2027}{18\textsuperscript{e} dimanche après la Pentecôte}
\jourcal{Dimanche}{19}{1}{vert}{\offgrand{18\textsuperscript{e} dimanche après la Pentecôte}}{dim.}{\comm{S. Janvier, Évêque, et ses Compagnons, Martyrs}\liturgie{Intr. \textit{Da pacem}}}
\jourcal{Lundi}{20}{0}{rouge}{\offpetit{S. Eustache et ses compagnons, Martyrs}}{dbl}{\comm{Vigile de S. Matthieu, Apôtre et Évangéliste}\liturgie{Intr. \textit{Sapiéntiam Sanctórum narrent populi}}\nov{Neuvaine à S. Michel — 1\textsuperscript{er} jour}{\pg{28}}}
\jourfete{Mardi}{21}{1}{rouge}{\offgrand{S. Matthieu, Apôtre et Évangéliste}}{2\textsuperscript{e} cl.}{\liturgie{Intr. \textit{Os justi meditábitur sapiéntiam}}\noticesaint{En Éthiopie, l’anniversaire de saint Matthieu, Apôtre et Évangéliste, qui, prêchant dans cette région, y souffrit le martyre. […]}\dev{Quinze mardis de S. Dominique — 15/15}{\pg{68}}\nov{Neuvaine à S. Michel — 2\textsuperscript{e} jour}{\pg{28}}}
\jourcal{Mercredi}{22}{0}{blanc}{\offpetit{S. Thomas de Villeneuve, Évêque et Confesseur}}{dbl}{\comm{S. Maurice et ses compagnons, Martyrs}\liturgie{Intr. \textit{Státuit ei Dóminus testaméntum pacis}}\nov{Neuvaine à S. Michel — 3\textsuperscript{e} jour}{\pg{28}}}
\semainedroite{du 19 au 25 septembre 2027}
\jourcal{Jeudi}{23}{0}{rouge}{\offpetit{S. Lin, Pape et Martyr}}{semi-dbl}{\comm{Ste Thècle, Vierge et Martyre}\liturgie{Intr. \textit{Sacerdótes Dei}}\nov{Neuvaine à S. Michel — 4\textsuperscript{e} jour}{\pg{28}}}
\jourcal{Vendredi}{24}{0}{blanc}{\offgrand{Notre-Dame de la Merci}}{dbl maj.}{\liturgie{Intr. \textit{Salve}}\jeune{Abstinence}\nov{Neuvaine à S. Michel — 5\textsuperscript{e} jour}{\pg{28}}}
\jourcal{Samedi}{25}{0}{blanc}{\offpetit{Office de la Sainte Vierge le samedi}}{simple}{\liturgie{Intr. \textit{Salve, sancta Parens}}\dev{Douze samedis avant l’Immaculée Conception — 2/12}{\pg{73}}\nov{Neuvaine à S. Michel — 6\textsuperscript{e} jour}{\pg{28}}}
\notes
\moisgauche{Octobre}{2027}{\devmois{Mois du saint Rosaire}{398}{0}}{Réjouissons-nous ensemble dans le Seigneur, car la fête que nous célébrons aujourd’hui est celle de la bienheureuse Vierge Marie. Cette solennité réjouit les Anges et tous en chœur louent le Fils de Dieu.}{Ps 44:2}{Saint Rosaire}
\moisdroite{Octobre 2027}{\rubmois{Fêtes principales}\lignemois{2}{Les Saints Anges gardiens}\lignemois{4}{S. François d’Assise, Confesseur}\lignemois{7}{Très Saint Rosaire de la Bienheureuse Vierge Marie}\lignemois{11}{Maternité de la Bienheureuse Vierge Marie}\lignemois{18}{S. Luc, Évangéliste}\lignemois{28}{Ss. Simon et Jude, Apôtres}\lignemois{31}{Notre-Seigneur Jésus-Christ Roi}\rubmois{Dévotions et neuvaines qui commencent}\lignemois{10}{Sept dimanches de l’Immaculée Conception\pg{65}}\lignemois{23}{Neuvaine de la Toussaint}\rubmois{Pénitence}\lignemoislong{Abstinence}{1, 8, 15, 22, 29}}
\semainegauche{du 26 septembre au 2 octobre 2027}{19\textsuperscript{e} dimanche après la Pentecôte}
\jourcal{Dimanche}{26}{1}{vert}{\offgrand{19\textsuperscript{e} dimanche après la Pentecôte}}{dim.}{\comm{S. Cyprien et Ste Justine, Martyrs}\liturgie{Intr. \textit{Salus pópuli ego sum}}\nov{Neuvaine à S. Michel — 7\textsuperscript{e} jour}{\pg{28}}}
\jourcal{Lundi}{27}{0}{rouge}{\offpetit{Ss. Côme et Damien, Martyrs}}{semi-dbl}{\liturgie{Intr. \textit{Sapiéntiam Sanctórum narrent populi}}\nov{Neuvaine à S. Michel — 8\textsuperscript{e} jour}{\pg{28}}}
\jourcal{Mardi}{28}{0}{rouge}{\offpetit{S. Wenceslas, Duc et Martyr}}{semi-dbl}{\liturgie{Intr. \textit{In virtúte tua}}\nov{Neuvaine à S. Michel — 9\textsuperscript{e} jour}{\pg{28}}\nov{Neuvaine du Rosaire — 1\textsuperscript{er} jour}{}}
\jourfete{Mercredi}{29}{1}{blanc}{\offgrand{Dédicace de S. Michel Archange}}{1\textsuperscript{re} cl.}{\liturgie{Intr. \textit{Benedícite Dóminum}}\nov{Neuvaine du Rosaire — 2\textsuperscript{e} jour}{}}
\semainedroite{du 26 septembre au 2 octobre 2027}
\jourcal{Jeudi}{30}{0}{blanc}{\offpetit{S. Jérôme, Confesseur et Docteur de l’Église}}{dbl}{\liturgie{Intr. \textit{In médio Ecclésiæ apéruit os…}}\dev{Heure sainte (veille du premier vendredi)}{\pg{69}}\nov{Neuvaine du Rosaire — 3\textsuperscript{e} jour}{}}
\jourcal{Vendredi}{1}{0}{blanc}{\offpetit{S. Remi, Évêque et Confesseur}}{simple}{\liturgie{Intr. \textit{Státuit ei Dóminus testaméntum pacis}}\jeune{Abstinence}\dev{Premier vendredi du mois}{\pg{69}}\nov{Neuvaine du Rosaire — 4\textsuperscript{e} jour}{}}
\jourcal{Samedi}{2}{0}{blanc}{\offgrand{Les Saints Anges gardiens}}{dbl maj.}{\liturgie{Intr. \textit{Benedícite Dóminum}}\noticesaint{La fête des saints Anges Gardiens.}\dev{Douze samedis avant l’Immaculée Conception — 3/12}{\pg{73}}\dev{Premier samedi du mois}{\pg{73}}\nov{Neuvaine du Rosaire — 5\textsuperscript{e} jour}{}}
\notes
\semainegauche{du 3 au 9 octobre 2027}{20\textsuperscript{e} dimanche après la Pentecôte}
\jourcal{Dimanche}{3}{1}{vert}{\offgrand{20\textsuperscript{e} dimanche après la Pentecôte}}{dim.}{\comm{Ste Thérèse de l’Enfant Jésus, Vierge}\liturgie{Intr. \textit{Omnia}}\nov{Neuvaine du Rosaire — 6\textsuperscript{e} jour}{}}
\jourcal{Lundi}{4}{0}{blanc}{\offgrand{S. François d’Assise, Confesseur}}{dbl maj.}{\liturgie{Intr. \textit{Mihi autem absit gloriári}}\noticesaint{À Assise, en Ombrie, l’anniversaire de saint François, Lévite et Confesseur, fondateur de trois Ordres : les Frères Mineurs, les Pauves Dames et les Frères et Sœurs de la Pénitence. […]}\nov{Neuvaine du Rosaire — 7\textsuperscript{e} jour}{}}
\jourcal{Mardi}{5}{0}{rouge}{\offpetit{S. Placide et ses Compagnons, Martyrs}}{simple}{\liturgie{Intr. \textit{Salus autem justórum a Dómino}}\nov{Neuvaine du Rosaire — 8\textsuperscript{e} jour}{}}
\jourcal{Mercredi}{6}{0}{blanc}{\offpetit{S. Bruno, Confesseur}}{dbl}{\liturgie{Intr. \textit{Os justi meditábitur sapiéntiam}}\dev{Premier mercredi : S. Joseph}{\pg{68}}\nov{Neuvaine du Rosaire — 9\textsuperscript{e} jour}{}}
\semainedroite{du 3 au 9 octobre 2027}
\jourfete{Jeudi}{7}{1}{blanc}{\offgrand{Très Saint Rosaire de la Bienheureuse Vierge Marie}}{2\textsuperscript{e} cl.}{\comm{S. Marc, Pape et Confesseur ; Ss. Serge, Bacchus, Marcel et Apulée, Martyrs}\liturgie{Intr. \textit{Gaudeámus omnes in Dómino}}\noticesaint{La fête de la bienheureuse Vierge Marie du Rosaire, et en même temps la commémoraison de Sainte Marie de la Victoire, fête que le souverain Pontife saint Pie V prescrivit de célébrer chaque année pour la […]}}
\jourcal{Vendredi}{8}{0}{blanc}{\offpetit{Ste Brigitte, Veuve}}{dbl}{\liturgie{Intr. \textit{Cognóvi}}\jeune{Abstinence}}
\jourcal{Samedi}{9}{0}{blanc}{\offpetit{S. Jean Léonardi, Confesseur}}{dbl}{\comm{Ss. Denis, Évêque, Rustique et Éleuthère, Martyrs}\liturgie{Intr. \textit{In sermónibus Dómini ópera ejus}}\dev{Douze samedis avant l’Immaculée Conception — 4/12}{\pg{73}}}
\notes
\semainegauche{du 10 au 16 octobre 2027}{21\textsuperscript{e} dimanche après la Pentecôte}
\jourcal{Dimanche}{10}{1}{vert}{\offgrand{21\textsuperscript{e} dimanche après la Pentecôte}}{dim.}{\comm{S. François de Borgia, Confesseur}\liturgie{Intr. \textit{In voluntáte tua}}\dev{Sept dimanches de l’Immaculée Conception — 1/7}{\pg{65}}}
\jourfete{Lundi}{11}{1}{blanc}{\offgrand{Maternité de la Bienheureuse Vierge Marie}}{2\textsuperscript{e} cl.}{\liturgie{Intr. \textit{Ecce Virgo concípiet}}\noticesaint{Fête de la Maternité de la bienheureuse Vierge Marie.}}
\jourcal{Mardi}{12}{0}{vert}{\offpetit{Férie}}{férie}{\liturgie{Intr. \textit{In voluntáte tua}}}
\jourcal{Mercredi}{13}{0}{blanc}{\offpetit{S. Édouard, Roi et Confesseur}}{semi-dbl}{\liturgie{Intr. \textit{Os justi meditábitur sapiéntiam}}}
\semainedroite{du 10 au 16 octobre 2027}
\jourcal{Jeudi}{14}{0}{rouge}{\offpetit{S. Calixte I\textsuperscript{er}, Pape et Martyr}}{dbl}{\liturgie{Intr. \textit{Státuit ei Dóminus testaméntum pacis}}}
\jourcal{Vendredi}{15}{0}{blanc}{\offpetit{Ste Thérèse, Vierge}}{dbl}{\liturgie{Intr. \textit{Dilexísti justítiam}}\jeune{Abstinence}}
\jourcal{Samedi}{16}{0}{blanc}{\offpetit{Ste Hedwige, Veuve}}{semi-dbl}{\liturgie{Intr. \textit{Cognóvi}}\dev{Douze samedis avant l’Immaculée Conception — 5/12}{\pg{73}}}
\notes
\semainegauche{du 17 au 23 octobre 2027}{22\textsuperscript{e} dimanche après la Pentecôte}
\jourcal{Dimanche}{17}{1}{vert}{\offgrand{22\textsuperscript{e} dimanche après la Pentecôte}}{dim.}{\comm{Ste Marguerite-Marie Alacoque, Vierge}\liturgie{Intr. \textit{Si iniquitátes observáveris}}\dev{Sept dimanches de l’Immaculée Conception — 2/7}{\pg{65}}}
\jourfete{Lundi}{18}{1}{rouge}{\offgrand{S. Luc, Évangéliste}}{2\textsuperscript{e} cl.}{\liturgie{Intr. \textit{Mihi autem nimis honoráti sunt…}}\noticesaint{En Bithynie, l’anniversaire du bienheureux Luc, Évangéliste. Après avoir beaucoup souffert pour le Nom du Christ, il mourut rempli du Saint-Esprit. […]}}
\jourcal{Mardi}{19}{0}{blanc}{\offpetit{S. Pierre d’Alcantara, Confesseur}}{dbl}{\liturgie{Intr. \textit{Justus ut palma florébit}}}
\jourcal{Mercredi}{20}{0}{blanc}{\offpetit{S. Jean de Kenty, Confesseur}}{dbl}{\liturgie{Intr. \textit{Miserátio hóminis circa próximum}}}
\semainedroite{du 17 au 23 octobre 2027}
\jourcal{Jeudi}{21}{0}{blanc}{\offpetit{S. Hilarion, Abbé}}{simple}{\comm{Ste Ursule et ses Compagnes, Vierges et Martyres}\liturgie{Intr. \textit{Os justi meditábitur sapiéntiam}}}
\jourcal{Vendredi}{22}{0}{vert}{\offpetit{Férie}}{férie}{\liturgie{Intr. \textit{Si iniquitátes observáveris}}\jeune{Abstinence}}
\jourcal{Samedi}{23}{0}{blanc}{\offpetit{Office de la Sainte Vierge le samedi}}{simple}{\liturgie{Intr. \textit{Salve, sancta Parens}}\dev{Douze samedis avant l’Immaculée Conception — 6/12}{\pg{73}}\nov{Neuvaine de la Toussaint — 1\textsuperscript{er} jour}{}}
\notes
\semainegauche{du 24 au 30 octobre 2027}{23\textsuperscript{e} dimanche après la Pentecôte}
\jourcal{Dimanche}{24}{1}{vert}{\offgrand{23\textsuperscript{e} dimanche après la Pentecôte}}{dim.}{\comm{S. Raphaël, Archange}\liturgie{Intr. \textit{Dicit Dóminus}}\dev{Sept dimanches de l’Immaculée Conception — 3/7}{\pg{65}}\nov{Neuvaine de la Toussaint — 2\textsuperscript{e} jour}{}}
\jourcal{Lundi}{25}{0}{rouge}{\offpetit{Ss. Chrysanthe et Darie, Martyrs}}{simple}{\liturgie{Intr. \textit{Intret in conspéctu tuo}}\nov{Neuvaine de la Toussaint — 3\textsuperscript{e} jour}{}}
\jourcal{Mardi}{26}{0}{rouge}{\offpetit{S. Évariste, Pape et Martyr}}{simple}{\liturgie{Intr. \textit{Státuit ei Dóminus testaméntum pacis}}\nov{Neuvaine de la Toussaint — 4\textsuperscript{e} jour}{}}
\jourcal{Mercredi}{27}{0}{violet}{\offpetit{Vigile des Ss. Simon et Jude, Apôtres}}{vigile}{\liturgie{Intr. \textit{Intret in conspéctu tuo}}\nov{Neuvaine de la Toussaint — 5\textsuperscript{e} jour}{}}
\semainedroite{du 24 au 30 octobre 2027}
\jourfete{Jeudi}{28}{1}{rouge}{\offgrand{Ss. Simon et Jude, Apôtres}}{2\textsuperscript{e} cl.}{\liturgie{Intr. \textit{Mihi autem nimis honoráti sunt…}}\noticesaint{En Perse, l’anniversaire des bienheureux Apôtres Simon le Cananéen, et Thaddée, appelé aussi Jude. Simon prêcha l’Évangile en Égypte, Thaddée en Mésopotamie. […]}\nov{Neuvaine de la Toussaint — 6\textsuperscript{e} jour}{}}
\jourcal{Vendredi}{29}{0}{vert}{\offpetit{Férie}}{férie}{\liturgie{Intr. \textit{Dicit Dóminus}}\jeune{Abstinence}\nov{Neuvaine de la Toussaint — 7\textsuperscript{e} jour}{}}
\jourcal{Samedi}{30}{0}{violet}{\offpetit{Vigile de la Toussaint (anticipée)}}{vigile}{\dev{Douze samedis avant l’Immaculée Conception — 7/12}{\pg{73}}\nov{Neuvaine de la Toussaint — 8\textsuperscript{e} jour}{}}
\notes
\moisgauche{Novembre}{2027}{\devmois{Mois des fidèles trépassés}{589}{1}}{Réjouissons-nous ensemble dans le Seigneur, car la fête que nous célébrons aujourd’hui est celle de tous les Saints. Cette solennité réjouit les Anges et tous en chœur louent le Fils de Dieu.}{Ps 32:1}{Toussaint}
\moisdroite{Novembre 2027}{\rubmois{Fêtes principales}\lignemois{1}{Toussaint}\lignemois{8}{Octave de la Toussaint}\lignemois{18}{Dédicace des Basiliques des Ss. Pierre et Paul}\rubmois{Pénitence}\lignemoislong{Abstinence}{5, 12, 19, 26}}
\semainegauche{du 31 octobre au 6 novembre 2027}{4\textsuperscript{e} dimanche après l’Épiphanie (repris après la Pentecôte)}
\jourfete{Dimanche}{31}{1}{blanc}{\offgrand{Notre-Seigneur Jésus-Christ Roi}}{1\textsuperscript{re} cl.}{\comm{4\textsuperscript{e} dimanche après l’Épiphanie (repris après la Pentecôte)}\liturgie{Intr. \textit{Dignus est Agnus}}\dev{Sept dimanches de l’Immaculée Conception — 4/7}{\pg{65}}\nov{Neuvaine de la Toussaint — 9\textsuperscript{e} jour}{}}
\jourfete{Lundi}{1}{1}{blanc}{\offgrand{Toussaint}}{1\textsuperscript{re} cl.}{\liturgie{Intr. \textit{Gaudeámus omnes in Dómino}}\noticesaint{La fête de tous les Saints. […]}}
\jourcal{Mardi}{2}{0}{noir}{\offpetit{Commémoraison de tous les fidèles défunts}}{dbl}{\liturgie{Intr. \textit{Réquiem ætérnam dona eis}}}
\jourcal{Mercredi}{3}{0}{blanc}{\offpetit{Troisième jour dans l’octave de la Toussaint}}{semi-dbl}{\liturgie{Intr. \textit{Gaudeámus omnes in Dómino}}\dev{Premier mercredi : S. Joseph}{\pg{68}}}
\semainedroite{du 31 octobre au 6 novembre 2027}
\jourcal{Jeudi}{4}{0}{blanc}{\offpetit{S. Charles Borromée, Évêque et Confesseur}}{dbl}{\comm{Quatrième jour dans l’octave de la Toussaint ; Ss. Vital et Agricola, Martyrs}\liturgie{Intr. \textit{Státuit ei Dóminus testaméntum pacis}}\dev{Heure sainte (veille du premier vendredi)}{\pg{69}}}
\jourcal{Vendredi}{5}{0}{blanc}{\offpetit{Cinquième jour dans l’octave de la Toussaint}}{semi-dbl}{\liturgie{Intr. \textit{Gaudeámus omnes in Dómino}}\jeune{Abstinence}\dev{Premier vendredi du mois}{\pg{69}}}
\jourcal{Samedi}{6}{0}{blanc}{\offpetit{Sixième jour dans l’octave de la Toussaint}}{semi-dbl}{\liturgie{Intr. \textit{Gaudeámus omnes in Dómino}}\dev{Douze samedis avant l’Immaculée Conception — 8/12}{\pg{73}}\dev{Premier samedi du mois}{\pg{73}}}
\notes
\semainegauche{du 7 au 13 novembre 2027}{5\textsuperscript{e} dimanche après l’Épiphanie (repris après la Pentecôte)}
\jourcal{Dimanche}{7}{1}{vert}{\offgrand{5\textsuperscript{e} dimanche après l’Épiphanie (repris après la Pentecôte)}}{dim.}{\comm{Septième jour dans l’octave de la Toussaint}\liturgie{Intr. \textit{Adoráte Deum}}\dev{Sept dimanches de l’Immaculée Conception — 5/7}{\pg{65}}}
\jourcal{Lundi}{8}{0}{blanc}{\offgrand{Octave de la Toussaint}}{dbl maj.}{\comm{Les quatre Saints couronnés, Martyrs}\liturgie{Intr. \textit{Gaudeámus omnes in Dómino}}}
\jourcal{Mardi}{9}{0}{blanc}{\offpetit{Dédicace de l’Archibasilique du Très-Saint-Sauveur}}{dbl}{\comm{S. Théodore, Martyr}\liturgie{Intr. \textit{Terríbilis est locus iste}}}
\jourcal{Mercredi}{10}{0}{blanc}{\offpetit{S. André Avellin, Confesseur}}{dbl}{\comm{Ss. Tryphon, Respice et Ste Nymphe, Vierge, Martyrs}\liturgie{Intr. \textit{Os justi meditábitur sapiéntiam}}}
\semainedroite{du 7 au 13 novembre 2027}
\jourcal{Jeudi}{11}{0}{blanc}{\offpetit{S. Martin, Évêque et Confesseur}}{dbl}{\comm{S. Menne, Martyr}\liturgie{Intr. \textit{Státuit ei Dóminus testaméntum pacis}}}
\jourcal{Vendredi}{12}{0}{rouge}{\offpetit{S. Martin I\textsuperscript{er}, Pape et Martyr}}{semi-dbl}{\liturgie{Intr. \textit{Sacerdótes Dei}}\jeune{Abstinence}}
\jourcal{Samedi}{13}{0}{blanc}{\offpetit{S. Didace, Confesseur}}{semi-dbl}{\liturgie{Intr. \textit{Justus ut palma florébit}}\dev{Douze samedis avant l’Immaculée Conception — 9/12}{\pg{73}}}
\notes
\semainegauche{du 14 au 20 novembre 2027}{6\textsuperscript{e} dimanche après l’Épiphanie (repris après la Pentecôte)}
\jourcal{Dimanche}{14}{1}{vert}{\offgrand{6\textsuperscript{e} dimanche après l’Épiphanie (repris après la Pentecôte)}}{dim.}{\comm{S. Josaphat, Évêque et Martyr}\liturgie{Intr. \textit{Adoráte Deum}}\dev{Sept dimanches de l’Immaculée Conception — 6/7}{\pg{65}}}
\jourcal{Lundi}{15}{0}{blanc}{\offpetit{S. Albert le Grand, Évêque, Confesseur et Docteur de l’Église}}{dbl}{\liturgie{Intr. \textit{In médio Ecclésiæ apéruit os…}}}
\jourcal{Mardi}{16}{0}{blanc}{\offpetit{Ste Gertrude, Vierge}}{dbl}{\liturgie{Intr. \textit{Dilexísti justítiam}}}
\jourcal{Mercredi}{17}{0}{blanc}{\offpetit{S. Grégoire le Thaumaturge, Évêque et Confesseur}}{semi-dbl}{\liturgie{Intr. \textit{Státuit ei Dóminus testaméntum pacis}}}
\semainedroite{du 14 au 20 novembre 2027}
\jourcal{Jeudi}{18}{0}{blanc}{\offgrand{Dédicace des Basiliques des Ss. Pierre et Paul}}{dbl maj.}{\liturgie{Intr. \textit{Terríbilis est locus iste}}\noticesaint{À Rome, la Dédicace des Basiliques de Saint Pierre et de Saint Paul, Apôtres. La première, ayant été reconstruite et agrandie, fut solennellement consacrée en ce jour par le souverain Pontife Urbain VIII. […]}}
\jourcal{Vendredi}{19}{0}{blanc}{\offpetit{Ste Élisabeth, Veuve}}{dbl}{\comm{S. Pontien, Pape et Martyr}\liturgie{Intr. \textit{Cognóvi}}\jeune{Abstinence}}
\jourcal{Samedi}{20}{0}{blanc}{\offpetit{S. Félix de Valois, Confesseur}}{dbl}{\liturgie{Intr. \textit{Justus ut palma florébit}}\dev{Douze samedis avant l’Immaculée Conception — 10/12}{\pg{73}}}
\notes
\semainegauche{du 21 au 27 novembre 2027}{24\textsuperscript{e} dimanche après la Pentecôte}
\jourcal{Dimanche}{21}{1}{vert}{\offgrand{24\textsuperscript{e} dimanche après la Pentecôte}}{dim.}{\comm{Présentation de la Bienheureuse Vierge Marie}\liturgie{Intr. \textit{Dicit Dóminus}}\dev{Sept dimanches de l’Immaculée Conception — 7/7}{\pg{65}}}
\jourcal{Lundi}{22}{0}{rouge}{\offpetit{Ste Cécile, Vierge et Martyre}}{dbl}{\liturgie{Intr. \textit{Loquébar de testimóniis tuis in…}}}
\jourcal{Mardi}{23}{0}{rouge}{\offpetit{S. Clément I\textsuperscript{er}, Pape et Martyr}}{dbl}{\comm{Ste Félicité, Martyre}\liturgie{Intr. \textit{Dicit Dóminus}}}
\jourcal{Mercredi}{24}{0}{rouge}{\offpetit{S. Jean de la Croix, Confesseur et Docteur de l’Église}}{dbl}{\comm{S. Chrysogone, Martyr}\liturgie{Intr. \textit{In médio Ecclésiæ apéruit os…}}}
\semainedroite{du 21 au 27 novembre 2027}
\jourcal{Jeudi}{25}{0}{rouge}{\offpetit{Ste Catherine d’Alexandrie, Vierge et Martyre}}{dbl}{\liturgie{Intr. \textit{Loquébar de testimóniis tuis in…}}}
\jourcal{Vendredi}{26}{0}{blanc}{\offpetit{S. Sylvestre, Abbé}}{dbl}{\comm{S. Pierre d’Alexandrie, Évêque et Martyr}\liturgie{Intr. \textit{Os justi meditábitur sapiéntiam}}\jeune{Abstinence}}
\jourcal{Samedi}{27}{0}{blanc}{\offpetit{Office de la Sainte Vierge le samedi}}{simple}{\liturgie{Intr. \textit{Salve, sancta Parens}}\dev{Douze samedis avant l’Immaculée Conception — 11/12 (le 12\textsuperscript{e} : 4 déc.)}{\pg{73}}}
\notes
